% The MATLAB/Octave chapter of the WFDB-x user's manual. It follows the sections of the
% Python chapter (chapters/python.tex), in the same order. Each \apientry of the section
% "The high-level API" is a function of matlab/+wfdbx listed under "High-level functions"
% or "Header metadata" in matlab/+wfdbx/Contents.m; its params name the positional
% arguments of the function line, then the name-value options in the order of the
% structure of defaults that the function parses (wfdbx.internal.hl_options). The
% examples of style "example" were run with GNU Octave 8.4, in a copy of tests/data that
% is the current folder (the paths /path/to/wfdb-x replaced by those of the repository);
% check_manual.py does not run them.
\chapter{MATLAB and Octave}
\label{chap:matlab}
\chapterlanguage{matlab}

The MATLAB/Octave implementation of WFDB-x is the package \code{wfdbx}: the folder
\file{matlab/+wfdbx} of the repository, with its internal functions in
\file{+wfdbx/+internal} and the command-line tools in \file{+wfdbx/+apps}. It is plain
MATLAB code, which runs unchanged in GNU Octave and in MATLAB, needs no toolbox, and never
calls the C library. Every function is called with the name of the package,
\code{wfdbx.NAME}.

It offers WFDB-x in three ways:
\begin{itemize}
\item \textbf{The high-level API} (\cref{sec:matlab-api}): functions that read and write
  whole records, headers, annotation files and header metadata as matrices of doubles,
  structures and cell arrays, in one call, with name-value options, and that throw MATLAB
  errors (with an identifier, such as \code{wfdbx:read\_record}) on failure. They read
  signal files in large blocks, and are the fastest way to read a record. They are an
  addition of WFDB-x: the C package has no equivalent. They are the recommended interface
  for new code and analysis.
\item \textbf{The WFDB library functions} under their C names, called with the name of
  the package as a prefix: \code{wfdbx.isigopen}, \code{wfdbx.getvec},
  \code{wfdbx.getann}, \code{wfdbx.putann}, and so on. They have the
  behavior of the C library (\cref{sec:matlab-library}). They are documented once, for
  all the languages, in \cref{chap:library}.
\item \textbf{The command-line tools} (\code{rdsamp}, \code{rdann}, \code{wrsamp},
  \code{xform} and 29 others), with the options, output and exit status of the C tools,
  run from the shell through an Octave script or a MATLAB shell script, or inside a
  session as functions (\cref{sec:matlab-tools}). They are documented in
  \cref{chap:tools}.
\end{itemize}
The high-level functions return exactly the samples and annotations that the library
functions (and the C library) give.

% ==========================================================================================
\section{Installation and requirements}
\label{sec:matlab-install}

\subsection{Requirements}

WFDB-x for MATLAB and Octave needs GNU Octave 6 or later, or MATLAB R2019b or later, and
no toolbox or other package. It uses only the language features and functions that both
products provide (no \code{string} objects, no \code{arguments} blocks, no Octave-only
syntax), so the same files run in both. It has no compiled code, and needs no C compiler,
no WFDB library and no libFLAC: the FLAC formats 508, 516 and 524 are read and written by
WFDB-x's own codec.

The code is tested with GNU Octave 8.4 (\cref{sec:matlab-testing}). It has not yet been
run in MATLAB: it is written to the rules that both products share, which are checked by
reading the code, and its MATLAB command-line driver, \file{matlab/bin/wfdbx\_matlab.sh},
has been checked in Octave only.

Remote files (\cref{chap:files}) are downloaded with the program \cmd{curl} when it is
installed, and otherwise with \code{webread} (Octave) or \code{websave} (MATLAB).

\subsection{Installing}

The package is used from a copy of the repository: its folder \file{matlab}, which holds
\file{+wfdbx}, is added to the path:

\begin{lstlisting}[language=Matlab]
addpath('/path/to/wfdb-x/matlab')
savepath      % optional: keep the folder on the path in later sessions
\end{lstlisting}

The folder \file{matlab} also holds copies of the license files of the repository:
\file{LICENSE} (the GNU Lesser General Public License, version 2.1, of the library),
\file{LICENSE-GPL} (the GNU General Public License, version 3, of the command-line tools:
\file{+wfdbx/+apps}, the drivers of \file{bin} and the dispatcher
\file{+wfdbx/+internal/cli\_main.m} and \file{app\_run.m}; the internal functions that
the library also uses, such as \file{tool\_command.m}, which \code{wfdbx.oannclose} calls
to run \cmd{sortann} in legacy mode, stay under the LGPL)
and \file{NOTICE}, which summarizes the licenses and gives the BSD-3-Clause notice of
libFLAC, whose stream decoder the FLAC code of \file{+wfdbx/+internal} translates. The
folder is the MATLAB and Octave distribution of WFDB-x: a copy of it that is passed on
keeps these three files.

\code{help wfdbx.Contents} lists the functions, by group (in MATLAB, \code{help wfdbx}
also does), and \code{help wfdbx.NAME} describes each one. \file{+wfdbx/Contents.m} gives
the version of the package, \code{Version 1.2.2}, and the command-line driver prints it:

\begin{lstlisting}[style=shell]
octave-cli --no-gui --norc -q matlab/bin/wfdbx_cli.m --version    # WFDB-x 1.2.2
\end{lstlisting}

\subsection{The sample record}
\label{sec:matlab-resources}

The repository includes the sample record \code{100s} in \file{tests/data}: the first
minute of record 100 of the MIT-BIH Arrhythmia Database, two ECG signals (\code{MLII} and
\code{V5}) sampled at 360\,Hz, with the reference beat annotations (annotator
\code{atr}), in the files \file{100s.hea}, \file{100s.dat} and \file{100s.atr}. The
package has no function that returns this folder: a program names it, or sets it as the
input directory. Work that uses the MIT-BIH Arrhythmia Database cites it as described on
PhysioNet.

% ==========================================================================================
\section{Quick start}
\label{sec:matlab-quickstart}

This example reads the header, ten seconds of the signals and the annotations of the
sample record \code{100s}, from the folder \file{tests/data} of a copy of the repository
(\code{'input\_dir'}), and computes the mean heart rate. It then writes the first ten
seconds as a new record, \code{first10}, in the current folder, with two metadata entries
in its header.

\begin{lstlisting}[language=Matlab,style=example]
addpath('/path/to/wfdb-x/matlab');
data = '/path/to/wfdb-x/tests/data';     % the folder of the sample record 100s

% The header only: no samples are read.
hdr = wfdbx.read_header('100s', 'input_dir', data);
fprintf('%d %g %d\n', hdr.nsig, hdr.fs, hdr.nsamples);   % 2 360 21600
fprintf('%s\n', hdr.siginfo.desc);                       % MLII, then V5

% Seconds 10 to 20 of both signals, in physical units: one row per sample, one column
% per signal.
rec = wfdbx.read_record('100s', 'sampfrom', '0:10', 'sampto', '0:20', ...
    'physical', true, 'input_dir', data);
disp(size(rec.signals));                % 3600 2
fprintf('%g %g\n', rec.signals(1, :));  % -0.39 -0.275

% The beat annotations, and the mean heart rate.
ann = wfdbx.read_annotations('100s', 'atr', 'input_dir', data);
fprintf('%d %s %s %s\n', numel(ann.samples), ann.symbols{1:3});   % 75 + N N
beats = ann.samples(wfdbx.isqrs(ann.codes) == 1);
rr = diff(beats) / ann.fs;              % RR intervals, in seconds
fprintf('%d beats, %.1f beats per minute\n', numel(beats), 60 / mean(rr));
% 74 beats, 73.9 beats per minute

% Write the first 10 seconds, in ADC units, as the record "first10" (first10.hea and
% first10.dat, in the current folder), with the signal specifications of 100s.
x = wfdbx.read_record('100s', 'sampto', '0:10', 'input_dir', data);
wfdbx.write_record('first10', x.signals, x.fs, 'fmt', 16, 'gains', x.gains, ...
    'baselines', x.baselines, 'adczeros', [x.siginfo.adczero], ...
    'adcres', [x.siginfo.adcres], 'units', 'mV', 'names', x.names, ...
    'metadata', {'Source', 'MIT-BIH record 100'});

% Header metadata: "# Key: value" lines of the header.
wfdbx.set_metadata('first10', 'Age', '69');
m = wfdbx.read_metadata('first10')';    % (transposed, for fprintf)
fprintf('%s: %s\n', m{:});      % Source: MIT-BIH record 100, then Age: 69
wfdbx.check_record('first10');  % prints: no problems found
\end{lstlisting}

The header \file{100s.hea} gives no units: the example writes \code{'mV'} explicitly,
since WFDB-x never assumes units, a gain, an ADC resolution or a sampling frequency that
a record does not give (\cref{chap:behavior}).

% ==========================================================================================
\section{The high-level API}
\label{sec:matlab-api}

The high-level API is made of the functions of this section, in the package
\code{wfdbx} (\code{wfdbx.read\_record}, \ldots), and of the structures of their results
(\cref{sec:matlab-results}). Each is one file of \file{matlab/+wfdbx};
\code{help wfdbx.NAME} gives its documentation.

\subsection{Conventions}
\label{sec:matlab-conventions}

\paragraph{Arguments and options.} The first arguments of a function are positional (the
record, the annotator, the samples, \ldots); the others are options, given as name-value
pairs after them, in any order: \code{wfdbx.read\_record('100s', 'sampto', 1000,
'physical', true)}. Option names are case-insensitive. An option that is not given keeps
its default, as listed in the call syntax of each entry; \code{[]} means ``not given'' for
the options whose default is \code{[]}. An unknown option, or a name without a value, is
an error (\cref{sec:matlab-warnings}). Text arguments are \code{char} row vectors or
MATLAB strings (converted to \code{char}).

\paragraph{Record names and directories.} A record is named as in the other parts of
WFDB-x (\cref{chap:records,chap:files}), without the \code{.hea} extension of its header
(the readers and the writers, \code{wfdbx.write\_record} and
\code{wfdbx.write\_annotations}, also accept a trailing \code{.hea}, which is removed). A
full path, absolute (\code{'/data/mitdb/100'}) or
in the home directory (\code{'\textasciitilde/Documents/100s'}), is used as given; a
relative path, such as \code{'100s'} or \code{'data/100s'}, is read from the input
directory if one is set, otherwise from the current folder, and written to the output
directory if one is set, otherwise to the current folder (missing folders are created);
a URL (\code{'https://physionet.org/files/mitdb/1.0.0/100'}) is read only if network
access is on, and cannot be written. A writer given a name with a path creates the files
in that folder, and the header names the record and its signal file without the path
(\code{'rec'}, \code{'rec.dat'}), so that the record opens under any path that leads to
it; in legacy mode, as in the original package, the header gives the name with its path.
The input and output directories are set with \code{wfdbx.wfdbx\_set\_input\_dir} and
\code{wfdbx.wfdbx\_set\_output\_dir} (\cref{sec:matlab-settings}) or with the
environment variables \code{WFDBX\_INPUT\_DIR} and \code{WFDBX\_OUTPUT\_DIR}. The
functions that read take the options
\code{'input\_dir'} and \code{'network'}, and those that write take \code{'dir'}: they
apply to that one call, and the previous settings are restored when it returns, also after
an error. \code{'input\_dir'} is optional, and applies to relative names only, not to full
paths or URLs; with \code{'dir'}, the record name must be a plain name (\code{'rec'}),
without a folder. The signal files that a header names are read from the header's
folder.

\paragraph{Times.} The options \code{'sampfrom'} and \code{'sampto'} are sample numbers
(whole numbers) or WFDB time strings, as \code{strtim} reads them: \code{'0:10'} (10
seconds), \code{'1:02:03.5'}, \code{'s3600'} (sample 3600), \code{'e'} (the end of the
record), or a time of day in brackets, \code{'[14:00:00]'}. A time in seconds needs the
sampling frequency; when the record does not give it, such a time is an error.

\paragraph{Numbers and arrays.} Samples, times, counts and codes are doubles. Signal
numbers start at 0, as in C. The signals of a record are a matrix with one row per sample
and one column per signal; the annotations of a file are column vectors and cell columns,
one row per annotation.

\paragraph{Per-signal and per-annotation options.} The options that describe each signal
(\code{'fmt'}, \code{'gains'}, \code{'units'}, \ldots) or each annotation
(\code{'subtypes'}, \code{'chans'}, \ldots) take one value for all, or a vector (for
numbers) or cell array (for text) with exactly one value per signal or per annotation.

\paragraph{Signal parameters.} No sampling frequency, gain, ADC resolution or units is
ever assumed. Where one is needed and the record does not give it, it must be given as an
argument; otherwise the function throws an error. A record whose header gives no
sampling frequency has \code{fs} 0.

\paragraph{Text.} Text is held as bytes: each character of a \code{char} vector is one byte
of a file (codes 0 to 255), so every byte sequence round-trips unchanged, and UTF-8 text
is stored as its bytes. In Octave, a \code{char} array holds bytes already; in MATLAB, a
character above 255 cannot be written to a header. The metadata functions also accept a
\code{uint8} vector of bytes as a key or value.

\paragraph{Legacy mode.} With \code{'legacy', true}, a function behaves as in the
original WFDB package, with its assumed signal parameters and its search for files
(\cref{chap:behavior}); the library's legacy mode is on for the duration of the call
only. With \code{false} (the default), the current setting of
\code{wfdbx.wfdbx\_set\_legacy} applies, which is off unless it was turned on.

\paragraph{State.} The functions that read and write records and annotation files call
\code{wfdbx.wfdbquit()}, which closes anything opened with the WFDB library functions,
and close their own files before they return, whatever the outcome.
\code{wfdbx.read\_record} sets the getvec mode for its own reading only. The functions
\code{wfdbx.read\_record}, \code{wfdbx.read\_header}, \code{wfdbx.read\_annotations},
\code{wfdbx.write\_record}, \code{wfdbx.write\_annotations} and
\code{wfdbx.check\_record} leave the getvec mode (\code{wfdbx.getgvmode}) and the time
resolution of output annotators (\code{wfdbx.getafreq}) as they were when they return,
also after an error: opening an EDF record sets both (it adds \code{WFDB\_HIGHRES} to the
getvec mode, and sets the time resolution to the record's sampling frequency when none is
set), as in C, and \code{wfdbx.wfdbquit()} resets neither. The metadata
functions open no signals, and leave the library's open files as they are.

\paragraph{The examples} of this section read the test records of WFDB-x (among them
\code{100s}) from the current folder, and write records of their own there: they run in a
copy of the folder \file{tests/data} of the repository, made the current folder (with
\file{matlab} on the path).

% ------------------------------------------------------------------------------------------
\subsection{Reading records}
\label{sec:matlab-reading}

\apientry{wfdbx.read\_record}{Read the signals of a record}
\begin{lstlisting}[language=Matlab,style=signature]
rec = wfdbx.read_record(record, 'channels', [], 'sampfrom', 0, 'sampto', [], ...
    'physical', false, 'highres', false, 'pad', false, 'input_dir', [], ...
    'network', [], 'legacy', false, 'checksum_warning', true)
\end{lstlisting}
Reads the signals of a record, or part of it, into a matrix of doubles with one row per
sample and one column per signal, together with the description of the record. The
samples are those that \code{getvec} returns (\cref{chap:library}), read in large blocks.
Invalid samples are \code{wfdbx.WFDB\_INVALID\_SAMPLE()} ($-32768$) in ADC units, and NaN
in physical units. When the read reaches the end of the record, the samples are compared
with the checksums of the header (\cref{sec:matlab-warnings}).
\begin{params}
\item[record] Record name (a \code{char} vector or a MATLAB string), without the \code{.hea} extension (a trailing \code{.hea}
  is also accepted): a full path, absolute (\code{'/data/mitdb/100'}) or in the home
  directory (\code{'\textasciitilde/Documents/100s'}), used as given; a relative path
  (\code{'100s'}, \code{'data/100s'}), read from the input directory if one is set,
  otherwise from the current folder; or a URL
  (\code{'https://physionet.org/files/mitdb/1.0.0/100'}), read only if network access is
  on (\cref{sec:record-paths}).
\item[channels] The signals to read, in the order in which they are returned: a vector of
  signal numbers (from 0), or a cell array of signal numbers and signal descriptions (as
  in \code{rec.names}, such as \code{'MLII'}), which may be mixed. Default \code{[]}: all
  the signals, in the order of the header.
\item[sampfrom] The first sample to read: a sample number (a whole number $\ge 0$) or a
  time string. Default \code{0}: the beginning of the record.
\item[sampto] The sample after the last one to read: a sample number (a whole number) or
  a time string. A value at or before \code{sampfrom} reads no samples. Default
  \code{[]}: the end of the record.
\item[physical] \code{false} (the default) to return ADC units; \code{true} to return
  physical units, computed as $(\mathit{adc} - \mathit{baseline}) / \mathit{gain}$. Every
  signal read must then have a gain.
\item[highres] \code{false} (the default) to read a multi-frequency record at its frame
  frequency, with the samples of the faster signals averaged to one per frame;
  \code{true} to read every sample, at the highest sampling frequency of the record (the
  samples of the slower signals are repeated). \code{rec.fs} is then that frequency.
\item[pad] \code{true} to replace each invalid sample by the previous valid sample of the
  same signal. Default \code{false}.
\item[input\_dir] The input directory for this call only (optional): a local folder, a
  \code{\textasciitilde} path or the URL of a directory, under which relative record names
  are read (\code{''}: the default). It does not apply to full paths or URLs. Default
  \code{[]}: the current setting (\code{wfdbx.wfdbx\_set\_input\_dir}).
\item[network] \code{true} to allow reading URLs for this call only, \code{false} to
  forbid it. Default \code{[]}: the current setting (\code{wfdbx.wfdbx\_set\_network}).
\item[legacy] \code{true} to behave as the original WFDB package (legacy mode, for this
  call only). Default \code{false}.
\item[checksum\_warning] \code{true} (the default) to issue the warning
  \code{wfdbx:checksum} when the samples do not match the checksums of the header;
  \code{false} for no warning. The samples are returned either way, and
  \code{rec.checksum\_ok} is \code{false}.
\end{params}
The memory allocated does not depend on the length that the header gives, which may be
far more than the signal files hold: the samples are read in chunks of about $2^{22}$,
joined at the end.
\returns{A structure (\cref{tab:matlab-record}).}
\raises{Identifier \code{wfdbx:read\_record}: the record cannot be read (``can't read
record REC''), or its signals cannot be opened; a signal file is shorter than its header
says, or cannot be read (a decoder or I/O error): ``can't read signal data of record REC
(getvec status -3)'', even if the checksums do not match either (the checksum error is
found in the frame in which the file ended), so that no incomplete frame is ever
returned; a requested signal is not in the record; \code{sampfrom} is negative, or
\code{sampfrom} or \code{sampto} is not a whole number (a fraction, NaN or an infinite
value: ``read\_record: sampfrom must be a whole number''), or is neither a number nor a
time string (``\ldots{} must be a sample number or a time string''); a time string needs the
sampling frequency and the header gives none; \code{'physical'} is \code{true} and a
signal read is uncalibrated; or an option is not valid. Warning \code{wfdbx:checksum} on
a checksum error (see \code{checksum\_warning}).}
\begin{lstlisting}[language=Matlab,style=example]
% Seconds 10 to 20 of lead MLII, in physical units.
rec = wfdbx.read_record('100s', 'channels', {'MLII'}, 'sampfrom', '0:10', ...
    'sampto', '0:20', 'physical', true);
disp(size(rec.signals));                         % 3600 1
fprintf('%s %g %d\n', rec.names{1}, rec.fs, rec.sampfrom);   % MLII 360 3600

% The first 3 samples of both signals, in ADC units, V5 first.
rec = wfdbx.read_record('100s', 'channels', [1 0], 'sampto', 3);
disp(rec.signals);          % three rows of 1011 995
disp(rec.names);            % {'V5', 'MLII'}
disp([rec.gains; rec.baselines]);      % 200 200; 1024 1024

% The whole record: the checksums are verified.
rec = wfdbx.read_record('100s');
fprintf('%d %d %d\n', size(rec.signals), rec.checksum_ok);    % 21600 2 1
\end{lstlisting}
A multi-frequency record (here the test record \code{mfreq}, whose signals have 1, 2
and 4 samples per frame at a frame frequency of 100\,Hz) is read at its frame frequency,
or at its highest sampling frequency:
\begin{lstlisting}[language=Matlab,style=example]
rec = wfdbx.read_record('mfreq');
fprintf('%g %d %d\n', rec.fs, size(rec.signals));     % 100 600 3
rec = wfdbx.read_record('mfreq', 'highres', true);
fprintf('%g %d %d\n', rec.fs, size(rec.signals));     % 400 2400 3
disp([rec.siginfo.spf]);                              % 1 2 4
\end{lstlisting}

\apientry{wfdbx.read\_header}{Read the header of a record}
\begin{lstlisting}[language=Matlab,style=signature]
hdr = wfdbx.read_header(record, 'input_dir', [], 'network', [], 'legacy', false)
\end{lstlisting}
Reads the header of a record, without opening its signal files. The result describes the
record and each of its signals, and holds the info strings and the metadata entries of the
header. For a multi-segment record, the master header is read, and \code{nsamples} is the
length of the whole record. The metadata are read from the \file{.hea} file only; an EDF
record (which has no \file{.hea} file) and a header read from the standard input
(\code{'-'}) have none.
\begin{params}
\item[record] Record name, without the \code{.hea} extension (a trailing \code{.hea}
  is also accepted): a full path, absolute (\code{'/data/mitdb/100'}) or in the home
  directory (\code{'\textasciitilde/Documents/100s'}), used as given; a relative path
  (\code{'100s'}, \code{'data/100s'}), read from the input directory if one is set,
  otherwise from the current folder; or a URL
  (\code{'https://physionet.org/files/mitdb/1.0.0/100'}), read only if network access is
  on (\cref{sec:record-paths}).
\item[input\_dir] The input directory for this call only (optional), under which
  relative record names are read; it does not apply to full paths or URLs (as for
  \code{wfdbx.read\_record}). Default \code{[]}: the current setting.
\item[network] \code{true} to allow reading URLs for this call only, \code{false} to
  forbid it. Default \code{[]}: the current setting.
\item[legacy] \code{true} to read the header as the original WFDB package does. Default
  \code{false}.
\end{params}
\returns{A structure (\cref{tab:matlab-header}).}
\raises{Identifier \code{wfdbx:read\_header}: the header cannot be found or read
(``can't read header for record REC''), or an option is not valid. Warning
\code{wfdbx:duplicate\_key} for a metadata key that the header has more than once.}
\begin{lstlisting}[language=Matlab,style=example]
hdr = wfdbx.read_header('100s');
fprintf('%s %g %d %d\n', hdr.name, hdr.fs, hdr.nsig, hdr.nsamples);  % 100s 360 2 21600
for k = 1:hdr.nsig
    s = hdr.siginfo(k);
    fprintf('%s %s %d %g %d\n', s.desc, s.fname, s.fmt, s.gain, s.adcres);
end
% MLII 100s.dat 212 200 11
% V5 100s.dat 212 200 11
disp(hdr.comments{1});          % " 69 M 1085 1629 x1": the first info string
disp(isempty(hdr.base_time));   % 1: the header gives no start time
\end{lstlisting}

\apientry{wfdbx.read\_annotations}{Read an annotation file}
\begin{lstlisting}[language=Matlab,style=signature]
ann = wfdbx.read_annotations(record, annotator, 'sampfrom', 0, 'sampto', [], ...
    'fs', [], 'input_dir', [], 'network', [], 'legacy', false)
ann = wfdbx.read_annotations(record, annotator, sampfrom, sampto, ...)
\end{lstlisting}
Reads the annotations of an annotation file, in file order, into column vectors and cell
columns. The sampling frequency, needed to convert times in seconds, is that of the
record's header, or, for a record without a header or whose header gives none, the option
\code{'fs'}. \code{'sampfrom'} and \code{'sampto'} may also be given as the third and
fourth arguments, without their names.
\begin{params}
\item[record] Record name, without the \code{.hea} extension (a trailing \code{.hea}
  is also accepted): a full path, absolute (\code{'/data/mitdb/100'}) or in the home
  directory (\code{'\textasciitilde/Documents/100s'}), used as given; a relative path
  (\code{'100s'}, \code{'data/100s'}), read from the input directory if one is set,
  otherwise from the current folder; or a URL
  (\code{'https://physionet.org/files/mitdb/1.0.0/100'}), read only if network access is
  on (\cref{sec:record-paths}). Its header, if there is one, gives the
  sampling frequency.
\item[annotator] Annotator name: the suffix of the annotation file. \code{'atr'} reads the
  file \file{RECORD.atr}.
\item[sampfrom] The first time to read: a sample number (a whole number $\ge 0$) or a
  time string. Default \code{0}: the beginning of the file.
\item[sampto] The time at which to stop: annotations at or after it are not read. A
  sample number (a whole number) or a time string. Default \code{[]}: the end of the
  file.
\item[fs] The sampling frequency, in Hz (a number $> 0$), to use when the record has no
  header or its header gives none; it is ignored otherwise. Default \code{[]}: none.
\item[input\_dir] The input directory for this call only (optional), under which
  relative record names are read; it does not apply to full paths or URLs (as for
  \code{wfdbx.read\_record}). Default \code{[]}: the current setting.
\item[network] \code{true} to allow reading URLs for this call only, \code{false} to
  forbid it. Default \code{[]}: the current setting.
\item[legacy] \code{true} to behave as the original WFDB package (which assumes 250\,Hz
  for a record without a frequency). Default \code{false}.
\end{params}
\returns{A structure (\cref{tab:matlab-annotations}). Its \code{fs} is the frequency in
which the times are counted: the one written in the annotation file, if there is one, or
else the sampling frequency; it is 0 if none of them is known.}
\raises{Identifier \code{wfdbx:read\_annotations}: the annotation file cannot be opened
(``can't read annotator ANN for record REC''), or it ends unexpectedly or cannot be read
before \code{sampto} (``can't read annotation file ANN of record REC (getann status
S)''); \code{'fs'} is given and is not $> 0$; \code{sampfrom} or \code{sampto} is not
valid, as for \code{wfdbx.read\_record} (``read\_annotations: sampfrom must be a whole
number'', ``\ldots{} must be a sample number or a time string'', ``\ldots{} is beyond
the range of a 64-bit integer'', ``read\_annotations: sampfrom must not be negative''),
which is checked before any file is opened; a time string needs the sampling frequency
and none is known; or an option is not valid. Whatever happens, the annotation file is
closed (\code{wfdbx.wfdbquit}) before the function returns.}
\begin{lstlisting}[language=Matlab,style=example]
ann = wfdbx.read_annotations('100s', 'atr', 'sampfrom', '0:10', 'sampto', '0:30');
fprintf('%d %g\n', numel(ann.samples), ann.fs);       % 24 360
fprintf('%d %s\n', ann.samples(1), ann.symbols{1});   % 3862 N

% The whole file: the first annotation is a rhythm change, with an aux string.
ann = wfdbx.read_annotations('100s', 'atr');
fprintf('%s %d %s\n', ann.symbols{1}, ann.codes(1) == wfdbx.RHYTHM(), ann.aux{1});
% + 1 (N
disp(unique(ann.symbols)');     % {'+', 'A', 'N'}
\end{lstlisting}
For a record without a header, \code{'fs'} gives the frequency in which times in seconds
are converted (the library reports on the standard error that it cannot open the
header):
\begin{lstlisting}[language=Matlab,style=example]
wfdbx.write_annotations('beats', 'qrs', [100 460 820], 'symbols', 'N');   % no header
ann = wfdbx.read_annotations('beats', 'qrs', 'sampto', '0:02', 'fs', 360);
fprintf('%d %d %g\n', ann.samples, ann.fs);           % 100 460 360
\end{lstlisting}

% ------------------------------------------------------------------------------------------
\subsection{Writing records}
\label{sec:matlab-writing}

\apientry{wfdbx.write\_record}{Write a header and a signal file}
\begin{lstlisting}[language=Matlab,style=signature]
wfdbx.write_record(record, signals, fs, 'fmt', 16, 'gains', [], 'baselines', [], ...
    'adczeros', 0, 'adcres', [], 'units', [], 'names', [], 'base_time', [], ...
    'comments', [], 'physical', false, 'samps_per_frame', [], 'metadata', [], ...
    'dir', [], 'legacy', false)
\end{lstlisting}
Writes signals as a new record: the header \file{RECORD.hea} and one signal file,
\file{RECORD.dat}, which holds all the signals (so they share one format). Existing files
of these names are replaced: both files are written to temporary files in the
destination folder (\file{.NAME.XXXXXXXX.tmp}, created with the permissions of the files
they replace), which replace the existing files only when both are complete, so that an
error leaves the existing record as it was, and no temporary file is left. The sampling
frequency, the gain and ADC resolution of
each signal, and the units of each calibrated signal must be given: there are no
defaults. The per-signal options (\code{'fmt'}, \code{'gains'}, \code{'baselines'},
\code{'adczeros'}, \code{'adcres'}, \code{'units'}, \code{'names'}) take one value for all
the signals, or a vector or cell array with one value per signal. All the arguments are
checked before any file is created or changed.
\begin{params}
\item[record] Record name (not empty), without the \code{.hea} extension (a trailing
  \code{.hea} is removed, as by the readers): a plain name
  (\code{'rec'}) or a relative path (\code{'data/rec'}), created in the output directory
  if one is set, otherwise in the current folder, missing folders included; or a full
  path, absolute or in the home directory (\code{'\textasciitilde/Documents/rec'}), used
  as given. With a path, the files are created in that folder, and the header names the
  record and its signal file without the path (\code{'rec'}, \code{'rec.dat'}), so that
  the record opens under any path that leads to it (in legacy mode, as in the original,
  the header gives the name with its path). A URL cannot be written.
\item[signals] The samples: a matrix (a 1-D or 2-D array) with one row per sample and one
  column per signal.
  In ADC units (the default), the values are integers within the range of a 32-bit
  integer; with \code{'physical', true}, they are in physical units.
\item[fs] The sampling frequency, in Hz: a finite number $> 0$.
\item[fmt] The storage format: 8, 16, 61, 80, 160, 212, 310, 311, 24, 32, or 508, 516 or
  524 (FLAC); one value, or one per signal, all equal (\cref{chap:formats}). Default
  \code{16}.
\item[gains] The gain of each signal, in ADC units per physical unit: a finite number;
  \code{0} marks an uncalibrated signal, which is allowed only for samples given in ADC
  units. Negative gains are allowed. Required (default \code{[]}: not given).
\item[baselines] The ADC value that corresponds to 0 physical units (a whole number within
  the range of a 32-bit integer; in a cell array, \code{[]} is the ADC zero). Default
  \code{[]}: the ADC zero of the signal.
\item[adczeros] The ADC value at the middle of the ADC range (a whole number within the
  range of a 32-bit integer; in a cell array, \code{[]} is 0). Default \code{0}.
\item[adcres] The ADC resolution, in bits (a whole number $\ge 1$, within the range of a
  32-bit integer). Required (default \code{[]}: not given).
\item[units] The physical units, such as \code{'mV'}: one string for all the signals, or a
  cell array with one per signal. Required for each signal whose gain is not 0. Default
  \code{[]}: not given.
\item[names] The signal descriptions, such as \code{'MLII'}: one string, or a cell array
  with one per signal. Default \code{[]}: no descriptions.
\item[base\_time] The start time, with an optional date: \code{'hh:mm:ss'} or
  \code{'hh:mm:ss dd/mm/yyyy'}, with optional decimals of the seconds, as
  \code{wfdbx.setbasetime} reads it. Default \code{[]}: no start time.
\item[comments] A cell array of info strings (or one string), written to the header as
  comment lines (\code{\#} followed by the string). Default \code{[]}: none.
\item[physical] \code{true} if \code{signals} are in physical units. They are converted to
  ADC units as $\mathrm{round}(x \cdot \mathit{gain}) + \mathit{baseline}$ (halves rounded
  away from zero); NaN becomes an invalid sample, and infinite values and values beyond
  the range of 32-bit ADC units are errors. Default \code{false}: ADC units.
\item[samps\_per\_frame] For a multi-frequency record, the number of samples per frame of
  each signal (integers $\ge 1$). \code{signals} then has one row per frame and
  \code{sum(samps\_per\_frame)} columns: the samples of signal 0, then those of signal 1,
  and so on. \code{fs} is then the frame frequency. Default \code{[]}: one sample per
  frame for every signal.
\item[metadata] Metadata entries (\cref{chap:metadata}) to write to the header after the
  comments, as \code{\# Key: value} (\code{\# Key:} for an empty value): an N-by-2 cell
  array with one entry per row, \code{\{KEY, VALUE; ...\}}. Keys and values follow the
  rules of \code{wfdbx.set\_metadata}; no two keys may be equal without regard to letter
  case, nor be the key of an entry in \code{'comments'}, and a value must not be
  \code{[]}. Default \code{[]}: none.
\item[dir] The folder (which must exist) in which to create the files, for this call
  only; \code{record} must then be a plain name, without a folder, and the record opens
  from that folder under its plain name. Default \code{[]}: the output directory.
\item[legacy] \code{true} to behave as the original WFDB package: the gain defaults to
  200, and the ADC resolution to one based on the format. Default \code{false}.
\end{params}
\returns{Nothing.}
\raises{Identifier \code{wfdbx:write\_record}: \code{record} is empty; \code{record}
has a character beyond 255 (``record name 'X' has a character that is not Latin-1 (one
byte per character)''; the same for a signal name, the units and a comment: ``signal
name 'X' \ldots'', ``units 'X' \ldots'', ``comment 'X' \ldots''; text is one byte per
character, so this can happen only in MATLAB, whose characters are 16-bit);
\code{signals} has more than two dimensions (``signals must be a 1-D or 2-D array''); a
\code{'samps\_per\_frame'} value is not an integer $\ge 1$ (``samps\_per\_frame must hold
integers >= 1''), or the number of columns of \code{signals} is not their sum (``signals
has N columns, but samps\_per\_frame adds up to M''); \code{signals} has no signals (``at
least one signal is required''); \code{fs} is not finite and $> 0$; a gain is not finite;
the gains, the ADC resolutions or the required units are not given; an ADC resolution, ADC
zero or baseline is not a whole number within the range of a 32-bit integer (text, NaN,
12.5: ``adcres must be whole numbers within the range of a 32-bit integer'', and the same
for \code{adczeros} and \code{baselines}); a per-signal option
does not have one value or one value per signal; a sample in ADC units is not an integer
or is outside the range of a 32-bit integer, or a sample in physical units is infinite or
its ADC value is outside that range; \code{'physical'} is \code{true} and a gain is 0;
\code{'base\_time'} cannot be parsed (``invalid base\_time 'T'''); the signals do not all
have the same format; an entry of \code{'metadata'} is not valid; \code{record} contains
a folder and \code{'dir'} is given, or \code{'dir'} does not exist; an option is not
valid; or the files cannot be created, or cannot be written completely, for example on a
full disk (``write error on record REC'', after the library's message; the existing files
are then left as they were).}
\begin{lstlisting}[language=Matlab,style=example]
% Two signals of random numbers, in millivolts, at 360 Hz: the record "rand"
% (rand.hea and rand.dat).
x = randn(3600, 2);
wfdbx.write_record('rand', x, 360, 'fmt', 16, 'physical', true, 'gains', 1000, ...
    'adcres', 16, 'units', 'mV', 'names', {'A', 'B'}, ...
    'base_time', '14:00:03 21/06/2024', 'comments', {'random numbers'}, ...
    'metadata', {'Generator', 'randn'});
rec = wfdbx.read_record('rand', 'physical', true);
fprintf('%d %d %s\n', size(rec.signals), rec.base_time);   % 3600 2 [14:00:03 21/06/2024]
disp(max(abs(rec.signals(:) - x(:))) <= 0.0005);           % 1: rounded to 1/1000 mV
disp(rec.comments');            % {'random numbers', ' Generator: randn'}

% An uncalibrated signal (gain 0, no units), in ADC units, in the folder "out".
mkdir('out');
wfdbx.write_record('counts', (0:999)', 100, 'gains', 0, 'adcres', 12, 'dir', 'out');
rec = wfdbx.read_record('out/counts');
disp(rec.signals(1:3)');        % 0 1 2
\end{lstlisting}
A multi-frequency record: signal \code{fast} has 2 samples per frame, signal
\code{slow} one; each row of \code{signals} is one frame, with the two samples of
\code{fast} and then the sample of \code{slow}:
\begin{lstlisting}[language=Matlab,style=example]
frames = 100;                                   % 1 second at 100 frames per second
fast = reshape(0:2 * frames - 1, 2, frames)';   % 200 samples per second
slow = repmat(500, frames, 1);                  % 100 samples per second
wfdbx.write_record('mf', [fast, slow], 100, 'gains', 200, 'adcres', 12, ...
    'units', 'mV', 'names', {'fast', 'slow'}, 'samps_per_frame', [2 1]);
rec = wfdbx.read_record('mf', 'highres', true);
fprintf('%g\n', rec.fs);        % 200
disp(rec.signals(1:3, :));      % 0 500; 1 500; 2 500
\end{lstlisting}

\apientry{wfdbx.write\_annotations}{Write an annotation file}
\begin{lstlisting}[language=Matlab,style=signature]
wfdbx.write_annotations(record, annotator, samples, 'symbols', [], 'codes', [], ...
    'subtypes', 0, 'chans', 0, 'nums', 0, 'aux', [], 'fs', [], 'dir', [], ...
    'sort', [], 'legacy', false)
\end{lstlisting}
Writes an annotation file, \file{RECORD.ANNOTATOR}, in the MIT format. An existing file of
that name is replaced: the new file is written to a temporary file in the same folder
(created with the permissions of the file it replaces), which replaces it only when it is
complete (and sorted), so that an error leaves it as it was. The record need not have a
header. The annotation types are
given as \code{'symbols'} or as \code{'codes'} (\cref{chap:formats}). Each per-annotation
option is one value for all the annotations, or a list with exactly one value per
annotation. If the annotations are not in time order, the file is sorted when it is
written, unless \code{'sort'} is \code{false}. All the arguments are checked before the
file is created or changed.
\begin{params}
\item[record] Record name, without the \code{.hea} extension (a trailing \code{.hea} is
  also accepted): a plain name or a relative path (\code{'data/100s'}) is created in the
  output directory if one is set, otherwise in the current folder; a full path, absolute
  or in the home directory (\code{'\textasciitilde/Documents/100s'}), is used as given.
  The file is \file{RECORD.ANNOTATOR} at that place (\file{data/100s.qrs}). A URL cannot
  be written.
\item[annotator] Annotator name: the suffix of the file, made of letters, digits,
  \code{\_}, \code{-} and \code{\textasciitilde}.
\item[samples] The annotation times, in samples: a vector (a 1-D array) of integers from
  $-2^{63}+1$ to $2^{63}-2$, one per annotation (the smallest and the largest 64-bit
  integers are reserved by the file format; doubles are exact up to $2^{53}$, and
  \code{int64} values are accepted).
\item[symbols] The mnemonics of the annotation types, such as \code{'N'} or \code{'V'}: a
  cell array with one per annotation, a \code{char} vector of one-character mnemonics
  (\code{'NNV'}), or one mnemonic for all. An unknown mnemonic becomes code 0
  (\code{NOTQRS}). Default \code{[]}: the types are given by \code{'codes'}.
\item[codes] The annotation codes (\code{wfdbx.NORMAL()}, \code{wfdbx.PVC()}, \ldots; 0
  to 58, as 59 to 63 are the control codes of the file format), instead of
  \code{'symbols'}: one per annotation, or one value for all. If both are given
  (as \code{wfdbx.read\_annotations} returns them), each symbol must be that of its code
  (an empty symbol is not checked), and the codes are written. Default \code{[]}.
\item[subtypes] The subtype of each annotation, stored as a signed byte (other values wrap
  around to $-128 \ldots 127$, as in C). Default \code{0}.
\item[chans] The channel number of each annotation, stored as an unsigned byte (other
  values wrap around to $0 \ldots 255$). Default \code{0}.
\item[nums] The annotator number of each annotation, stored as a signed byte. Default
  \code{0}.
\item[aux] The auxiliary text of each annotation: a cell array with one per annotation,
  with \code{[]} where there is none, or one \code{char} vector for all. Default
  \code{[]}: none.
\item[fs] If given (a finite number $> 0$), the frequency in which the times are
  counted, written to the file. Default \code{[]}: none is written.
\item[dir] The folder (which must exist) in which to create the file, for this call
  only; \code{record} must then be a plain name, without a folder: the file is
  \file{DIR/RECORD.ANNOTATOR}. Default \code{[]}: the output directory.
\item[sort] \code{true} to always sort the file, \code{false} to keep the given order
  (without the library's advice to run \code{sortann}), for this call only. Default
  \code{[]}: the environment variable \code{WFDBANNSORT} decides (sorting is on unless it
  is \code{0}).
\item[legacy] \code{true} to behave as the original WFDB package. Default \code{false}.
\end{params}
\returns{Nothing.}
\raises{Identifier \code{wfdbx:write\_annotations}: \code{samples} is a matrix or an
array of more dimensions (``samples must be a 1-D array''; the same for \code{'codes'},
\code{'subtypes'}, \code{'chans'}, \code{'nums'}, and for \code{'symbols'} given with
\code{'codes'}); neither \code{'symbols'} nor \code{'codes'} is given for one annotation
or more (``symbols or codes must be given''; with no annotations, an empty file is
written); a per-annotation
option given as a list does not have one value per annotation (``\ldots{} must have one
value per annotation''); a time is not an integer within the range of a 64-bit integer
(NaN, an infinite value or a fraction); a code, subtype, channel or annotator number is
NaN or infinite; \code{'fs'} is not finite and $> 0$; a time is $-2^{63}$ or, as an
\code{int64}, $2^{63}-1$ (``annotation time T cannot be written (annotation times must be
integers from -9223372036854775807 to 9223372036854775806)''); a code is not from 0 to 58
(``annotation code C is reserved (59 to 63 are the control codes of the annotation file
format)'', ``annotation code C is out of range (0 to 58)''); both \code{'symbols'} and
\code{'codes'} are given and disagree (``symbols and codes disagree at annotation K'', K
counting from 0); an aux text has a character beyond 255 (``aux text 'X' has a character
that is not Latin-1 (one byte per character)'', in MATLAB only); \code{record} contains a
folder and \code{'dir'} is given, or \code{'dir'} does not exist; an option is not valid;
or the file cannot be created, or cannot be written completely, or the library reports an
error while writing it (``write error on annotator ANN for record REC''; the existing file
is then left as it was).}
\begin{lstlisting}[language=Matlab,style=example]
% Three beats and a rhythm change, by symbol, with an aux string for the rhythm.
wfdbx.write_annotations('100s', 'test', [18 77 370 662], ...
    'symbols', {'+', 'N', 'V', 'N'}, 'aux', {'(N', [], [], []});
ann = wfdbx.read_annotations('100s', 'test');
fprintf('%s', ann.symbols{:}); fprintf(' %s\n', ann.aux{1});    % +NVN (N

% A copy of the reference annotations, as read_annotations returns them.
a = wfdbx.read_annotations('100s', 'atr');
wfdbx.write_annotations('100s', 'copy', a.samples, 'symbols', a.symbols, ...
    'codes', a.codes, 'subtypes', a.subtypes, 'chans', a.chans, 'nums', a.nums, ...
    'aux', a.aux);
b = wfdbx.read_annotations('100s', 'copy');
disp([isequal(b.samples, a.samples), isequal(b.symbols, a.symbols), ...
    isequal(b.aux, a.aux)]);    % 1 1 1

% Annotations out of time order are sorted when the file is written (the library
% reports it on the standard error).
wfdbx.write_annotations('100s', 'unsorted', [500 100], 'codes', wfdbx.NORMAL());
ann = wfdbx.read_annotations('100s', 'unsorted');
disp(ann.samples');             % 100 500
\end{lstlisting}

% ------------------------------------------------------------------------------------------
\subsection{Checking records}
\label{sec:matlab-checking}

\apientry{wfdbx.check\_record}{Check a record and its annotation files}
\begin{lstlisting}[language=Matlab,style=signature]
report = wfdbx.check_record(record, 'annotators', {}, 'input_dir', [], 'network', [])
\end{lstlisting}
Checks the header of a record, its signal specifications, every sample of its signal
files, each segment of a multi-segment record, and the annotation files named, and
reports each problem found. A problem in the record never throws an error: it is
reported, with a severity (\code{'error'}: the record is damaged or does not follow the
WFDB format, and programs will read wrong data or fail; \code{'warning'}: the record can
be read, but something is missing or unusual, and may give wrong results; \code{'info'}:
a fact worth knowing, such as invalid samples). The problems, their codes and their
messages are the same in the seven languages (\cref{tab:python-issue-codes}). The checks
run in normal mode, even if legacy mode is on, and the library prints no messages during
the call. They are similar to those of Benjamin Moody's \code{wfdbcheck}. Called without
an output argument, the function displays the lines of the report.
\begin{params}
\item[record] Record name, without the \code{.hea} extension (a trailing \code{.hea}
  is also accepted): a full path, absolute (\code{'/data/mitdb/100'}) or in the home
  directory (\code{'\textasciitilde/Documents/100s'}), used as given; a relative path
  (\code{'100s'}, \code{'data/100s'}), read from the input directory if one is set,
  otherwise from the current folder; or a URL
  (\code{'https://physionet.org/files/mitdb/1.0.0/100'}), read only if network access is
  on (\cref{sec:record-paths}).
\item[annotators] The annotator names whose files are checked: a cell array, such as
  \code{\{'atr'\}}, or one name as a \code{char} vector. Default \code{\{\}}: none.
\item[input\_dir] The input directory for this call only (optional), under which
  relative record names are read; it does not apply to full paths or URLs (as for
  \code{wfdbx.read\_record}). Default \code{[]}: the current setting.
\item[network] \code{true} to allow reading URLs for this call only, \code{false} to
  forbid it. Default \code{[]}: the current setting.
\end{params}
\returns{A structure (\cref{tab:matlab-checkreport}): \code{report.ok} is \code{true} if
no problem is an error, \code{report.issues} is a structure array of the problems
(\cref{tab:matlab-issue}), and \code{report.lines} is a cell column with one line per
problem, or the one line \code{no problems found}.}
\raises{Identifier \code{wfdbx:check\_record}, only for invalid arguments:
\code{record} is not a string or is empty, an annotator name is not a string, or an
option is not valid. A record that cannot be read is reported (code
\code{header\_unreadable}), not thrown.}
\begin{lstlisting}[language=Matlab,style=example]
report = wfdbx.check_record('100s', 'annotators', {'atr'});
disp(report.ok);                % 1: no errors
fprintf('%s\n', report.lines{:});
% warning: signal 0: the units are not specified
% warning: signal 1: the units are not specified

% A damaged record: one sample of a 12-bit signal is overwritten with 2048.
wfdbx.write_record('damaged', (0:99)', 100, 'gains', 100, 'adcres', 12, ...
    'units', 'mV', 'names', 'ramp');
fid = fopen('damaged.dat', 'r+b');
fseek(fid, 20, 'bof');          % sample 10 (2 bytes per sample in format 16)
fwrite(fid, [0 8], 'uint8');    % 0x0800 = 2048, little-endian
fclose(fid);
report = wfdbx.check_record('damaged');
disp(report.ok);                % 0
for k = 1:numel(report.issues)
    i = report.issues(k);
    fprintf('%s %s %s\n', i.severity, i.code, i.where);
end
% error sample_range signal 0
% error checksum signal 0
disp(report.lines{1});
% error: signal 0: 1 sample is outside the ADC range (-2048 to 2047): sample 10 (value 2048)
\end{lstlisting}

% ------------------------------------------------------------------------------------------
\subsection{Header metadata}
\label{sec:matlab-metadata}

These functions read and change the metadata of a header: its info strings of the form
\code{\# Key: value}, such as \code{\# Age: 63} (\cref{chap:metadata}). In short: keys are
compared without regard to letter case (A--Z and a--z only); a key must not be empty and
must not contain \code{:}, \code{\#}, a line break or a NUL character, and a value must
not contain a line break or a NUL character; spaces, tabs, VT, FF and CR at either end of
keys and values are removed; text is one byte per character, and keys and values may be
\code{char} vectors, MATLAB strings or \code{uint8} vectors of bytes. The entries are an
N-by-2 cell array of \code{char} row vectors, one entry per row, \code{\{KEY, VALUE;
...\}}, in file order (\code{cell(0, 2)} if there is none); a structure cannot hold keys
with spaces, such as \code{Date of birth}, so none is used. The info strings are the
lines that begin with \code{\#} after the record line, which is found as the header
parser of the library finds it, in every mode: the first line that holds something
other than white space (spaces, tabs, CRs and LFs, up to a NUL character) and does not
begin, after white space, with \code{\#}. Comment lines, and empty or blank lines,
before the record line are never read or changed. A change rewrites the header
in one step, through a temporary file, keeping every other byte of it; all the arguments
are checked first, so an error leaves the header unchanged. Only the \file{.hea} file is
read (not a \file{.info} file); for a multi-segment record, it is the master header,
unless a segment is named. Metadata often holds information about patients: WFDB-x stores
it as plain text, exactly as given, and does not de-identify it.

Each function has two error identifiers (\cref{sec:matlab-warnings}): \code{wfdbx:NAME}
for an invalid argument, and \code{wfdbx:NAME:file} for a header that cannot be read or
written, or a \code{'dir'} that does not exist (NAME is the name of the function, such as
\code{set\_metadata}).

\begin{phibox}
A header can hold identifying data about a patient (a name, a date of birth, a medical
record number). The metadata functions read and write such data as given. Records that
are shared must be de-identified first.
\end{phibox}

\begin{note}[Permissions of a changed header]
The temporary file is renamed over the header (\code{rename} in Octave, \code{movefile}
in MATLAB), so the header keeps its permissions only if the temporary file has them. The
temporary file is created with the permissions of the header, whatever the name of the
header (parentheses, spaces and any other characters included), and no shell command is
run on the name: Octave, which has no \code{chmod} function, creates it under a
\code{umask} that leaves exactly the read and write permissions of the header, so that a
private (0600) header stays private, and is never readable by others while it is
rewritten; MATLAB gives it the read, write and execute permissions of the header with
\code{fileattrib}. In Octave, execute permissions (which a header does not normally have)
are not kept. The files that \code{wfdbx.write\_record} and \code{wfdbx.write\_annotations}
replace keep their permissions in the same way.
\end{note}

\apientry{wfdbx.read\_metadata}{Read the metadata of a record}
\begin{lstlisting}[language=Matlab,style=signature]
M = wfdbx.read_metadata(record, 'input_dir', [], 'network', [])
\end{lstlisting}
Reads the metadata entries of the header of a record, in file order. An info string is an
entry if it contains a colon: the key is the text before the first colon, the value the
rest of the line, both without spaces at either end. Info strings without a colon, and
those whose key is empty or contains \code{\#}, are free comments, which are not returned.
\begin{params}
\item[record] Record name, without the \code{.hea} extension (a trailing \code{.hea}
  is also accepted): a full path, absolute (\code{'/data/mitdb/100'}) or in the home
  directory (\code{'\textasciitilde/Documents/100s'}), used as given; a relative path
  (\code{'100s'}, \code{'data/100s'}), read from the input directory if one is set,
  otherwise from the current folder; or a URL
  (\code{'https://physionet.org/files/mitdb/1.0.0/100'}), read only if network access is
  on (\cref{sec:record-paths}).
\item[input\_dir] The input directory for this call only (optional), under which
  relative record names are read; it does not apply to full paths or URLs (as for
  \code{wfdbx.read\_record}). Default \code{[]}: the current setting.
\item[network] \code{true} to allow reading URLs for this call only, \code{false} to
  forbid it. Default \code{[]}: the current setting.
\end{params}
\returns{An N-by-2 cell array of \code{char} row vectors, \code{\{KEY, VALUE; ...\}}, in
file order. A key that appears more than once (ignoring letter case) is returned once,
with its first value, and the warning \code{wfdbx:duplicate\_key} is issued
(\cref{sec:matlab-warnings}).}
\raises{Identifier \code{wfdbx:read\_metadata}: \code{record} is empty or \code{'-'} (the
standard input), or an option is not valid. Identifier \code{wfdbx:read\_metadata:file}:
the header cannot be found or read (``can't read header for record REC'').}
\begin{lstlisting}[language=Matlab,style=example]
wfdbx.write_record('pat', zeros(360, 1), 360, 'gains', 200, 'adcres', 12, ...
    'units', 'mV', 'names', 'ECG', 'comments', {'recorded at rest'}, ...
    'metadata', {'Age', '63'; 'Sex', 'F'});
M = wfdbx.read_metadata('pat');
disp(size(M));                  % 2 2
M = M';
fprintf('%s: %s\n', M{:});      % Age: 63, then Sex: F
hdr = wfdbx.read_header('pat');
disp(hdr.comments');
% {'recorded at rest', ' Age: 63', ' Sex: F'}: the free comment is not an entry
\end{lstlisting}

\apientry{wfdbx.get\_metadata}{Return the value of a metadata entry}
\begin{lstlisting}[language=Matlab,style=signature]
value = wfdbx.get_metadata(record, key, 'input_dir', [], 'network', [])
\end{lstlisting}
Returns the value of one metadata entry of the header of a record.
\begin{params}
\item[record] Record name (as for \code{wfdbx.read\_metadata}).
\item[key] The key, such as \code{'Age'}. Letter case is ignored, and spaces at either end
  are removed.
\item[input\_dir] The input directory for this call only (optional), under which a
  relative record name is read (as for \code{wfdbx.read\_metadata}). Default \code{[]}:
  the current setting.
\item[network] \code{true} to allow reading URLs for this call only, \code{false} to
  forbid it. Default \code{[]}: the current setting.
\end{params}
\returns{The value, a \code{char} row vector (empty for an entry with an empty value), or
\code{[]} if the header has no entry for \code{key} (\code{ischar} tells them apart). If
it has more than one, the first is returned, and the warning \code{wfdbx:duplicate\_key}
is issued.}
\raises{Identifier \code{wfdbx:get\_metadata}: \code{record} is empty or \code{'-'},
\code{key} is missing, not a string (or \code{uint8} vector) or not a valid key, or an
option is not valid. Identifier \code{wfdbx:get\_metadata:file}: the header cannot be
found or read.}
\begin{lstlisting}[language=Matlab,style=example]
wfdbx.write_record('pat', zeros(360, 1), 360, 'gains', 200, 'adcres', 12, ...
    'units', 'mV', 'metadata', {'Age', '63'; 'Sex', 'F'});
disp(wfdbx.get_metadata('pat', 'age'));         % 63
v = wfdbx.get_metadata('pat', 'Height');
disp([isempty(v), ischar(v)]);                  % 1 0: there is no such entry
\end{lstlisting}

\apientry{wfdbx.set\_metadata}{Set the value of a metadata entry}
\begin{lstlisting}[language=Matlab,style=signature]
wfdbx.set_metadata(record, key, value, 'input_dir', [], 'network', [], 'dir', [])
\end{lstlisting}
Sets the value of a metadata entry in the header of a record. If the header has an entry
for \code{key} (ignoring letter case), its line is replaced by \code{\# Key: value}, at
the same position and with the same line ending, with the spelling of \code{key} given
here (even if the value is the same); otherwise the entry is added after the last info
string of the header (or, if it has none, after its last line that is not blank), with
the line ending of the record line. The entry is written as \code{\# Key: value}, or
\code{\# Key:} if the value is empty. Every other line is kept, byte for byte, including
its line ending. The header is changed where it was read (not in the output directory),
unless \code{'dir'} is given.
\begin{params}
\item[record] Record name (as for \code{wfdbx.read\_metadata}). The header is changed
  where it was read, unless \code{'dir'} is given; a header read from a URL can only be
  written elsewhere, with \code{'dir'}.
\item[key] The key, such as \code{'Age'} or \code{'Date of birth'}. Spaces at either end
  are removed. It must not be empty, and must not contain \code{:}, \code{\#}, a line
  break or a NUL character.
\item[value] The value: text, which may be a sentence. Spaces at either end are removed.
  It must not contain a line break or a NUL character; it may be empty (\code{''}).
\item[input\_dir] The input directory in which the header is read, for this call only.
  Default \code{[]}: the current setting.
\item[network] \code{true} to allow reading the header from a URL, for this call only,
  \code{false} to forbid it. Default \code{[]}: the current setting.
\item[dir] A folder, which must exist, in which to write the changed header, under the
  same file name, instead of changing the header itself (which is then unchanged); the
  signal files are not copied. An unchanged header is also written there. Default
  \code{[]} (or \code{''}): the header itself is changed.
\end{params}
\returns{Nothing.}
\raises{Identifier \code{wfdbx:set\_metadata}: \code{key} or \code{value} is missing, not
a string (or \code{uint8} vector) or not valid; the header has more than one entry for
\code{key} (``\ldots{} remove them first''); it has no record line; it was read from a URL
and \code{'dir'} is not given; \code{record} is empty or \code{'-'}; \code{'dir'} is not a
string; or an option is not valid. Identifier \code{wfdbx:set\_metadata:file}: the header
cannot be read, or the changed header cannot be written (``can't write header for record
REC''); or \code{'dir'} does not exist (``directory DIR does not exist'').}
\begin{lstlisting}[language=Matlab,style=example]
wfdbx.write_record('pat', zeros(360, 1), 360, 'gains', 200, 'adcres', 12, ...
    'units', 'mV', 'metadata', {'Age', '63'; 'Sex', 'F'});
wfdbx.set_metadata('pat', 'AGE', '64');                 % replaces "# Age: 63"
wfdbx.set_metadata('pat', 'Diagnosis', 'sinus rhythm, no ectopy');   % a new entry
M = wfdbx.read_metadata('pat')';
fprintf('%s: %s\n', M{:});
% AGE: 64
% Sex: F
% Diagnosis: sinus rhythm, no ectopy

% Write the changed header to another folder; pat.hea is left unchanged.
mkdir('changed');
wfdbx.set_metadata('pat', 'Sex', 'M', 'dir', 'changed');
fprintf('%s %s\n', wfdbx.get_metadata('pat', 'Sex'), ...
    wfdbx.get_metadata('changed/pat', 'Sex'));         % F M
\end{lstlisting}

\apientry{wfdbx.remove\_metadata}{Remove a metadata entry}
\begin{lstlisting}[language=Matlab,style=signature]
n = wfdbx.remove_metadata(record, key, 'input_dir', [], 'network', [], 'dir', [])
\end{lstlisting}
Removes every entry for \code{key} (ignoring letter case) from the header of a record,
line ending included, so it also repairs a header that has more than one. Other lines are
kept, byte for byte. Nothing is written if the header has no entry for \code{key}, unless
\code{'dir'} is given.
\begin{params}
\item[record] Record name (as for \code{wfdbx.set\_metadata}).
\item[key] The key (as for \code{wfdbx.set\_metadata}).
\item[input\_dir] The input directory in which the header is read, for this call only.
  Default \code{[]}: the current setting.
\item[network] \code{true} to allow reading the header from a URL, for this call only,
  \code{false} to forbid it. Default \code{[]}: the current setting.
\item[dir] A folder, which must exist, in which to write the changed header instead of
  changing the header itself. Default \code{[]}: the header itself is changed.
\end{params}
\returns{The number of lines removed (0 if the header has no entry for \code{key}).}
\raises{As \code{wfdbx.set\_metadata}, with the identifiers
\code{wfdbx:remove\_metadata} and \code{wfdbx:remove\_metadata:file}, except that more
than one entry for \code{key} is not an error.}
\begin{lstlisting}[language=Matlab,style=example]
wfdbx.write_record('pat', zeros(360, 1), 360, 'gains', 200, 'adcres', 12, ...
    'units', 'mV', 'metadata', {'Age', '63'; 'Sex', 'F'});
disp(wfdbx.remove_metadata('pat', 'sex'));      % 1
disp(wfdbx.remove_metadata('pat', 'Sex'));      % 0: there is no such entry now
disp(wfdbx.read_metadata('pat'));               % {'Age', '63'}
\end{lstlisting}

\apientry{wfdbx.update\_metadata}{Make several changes to the metadata}
\begin{lstlisting}[language=Matlab,style=signature]
wfdbx.update_metadata(record, changes, 'input_dir', [], 'network', [], 'dir', [])
\end{lstlisting}
Makes several changes to the metadata of a record, in one write of the header. Each
change is applied in order, as \code{wfdbx.set\_metadata} does, or, for a value of
\code{[]}, as \code{wfdbx.remove\_metadata} does. All the changes are checked first: if one
is not valid, none is made.
\begin{params}
\item[record] Record name (as for \code{wfdbx.set\_metadata}).
\item[changes] An N-by-2 cell array with one change per row, \code{\{KEY, VALUE; ...\}};
  a value of \code{[]} (an empty double) removes the entry, and \code{''} sets an empty
  value. No two keys may be equal without regard to letter case. An empty cell array
  (\code{\{\}} or \code{cell(0, 2)}) changes nothing.
\item[input\_dir] The input directory in which the header is read, for this call only.
  Default \code{[]}: the current setting.
\item[network] \code{true} to allow reading the header from a URL, for this call only,
  \code{false} to forbid it. Default \code{[]}: the current setting.
\item[dir] A folder, which must exist, in which to write the changed header instead of
  changing the header itself. Default \code{[]}: the header itself is changed.
\end{params}
\returns{Nothing.}
\raises{As \code{wfdbx.set\_metadata}, with the identifiers
\code{wfdbx:update\_metadata} and \code{wfdbx:update\_metadata:file}; also an argument
error if \code{changes} is not an N-by-2 cell array, or if two of its keys differ only in
letter case.}
\begin{lstlisting}[language=Matlab,style=example]
wfdbx.write_record('pat', zeros(360, 1), 360, 'gains', 200, 'adcres', 12, ...
    'units', 'mV', 'metadata', {'Age', '63'; 'Sex', 'F'});
wfdbx.update_metadata('pat', {'Age', '64'; 'Sex', []; 'Dx', '164889003'});
M = wfdbx.read_metadata('pat')';
fprintf('%s: %s\n', M{:});      % Age: 64, then Dx: 164889003
try
    wfdbx.update_metadata('pat', {'Dx', '426783006'; 'dx', []});
catch err
    disp(err.identifier);       % wfdbx:update_metadata
    disp(err.message);          % metadata key 'dx' is given more than once (...)
end
disp(wfdbx.get_metadata('pat', 'Dx'));          % 164889003: nothing was changed
\end{lstlisting}

% ------------------------------------------------------------------------------------------
\subsection{The result structures}
\label{sec:matlab-results}

The readers return structures, whose fields are listed in
\cref{tab:matlab-record,tab:matlab-header,tab:matlab-annotations};
\code{wfdbx.check\_record} returns a structure that holds a structure array of issues
(\cref{tab:matlab-checkreport,tab:matlab-issue}). The signals of a record are described by
\code{wfdbx.siginfo} structures (\cref{tab:matlab-siginfo}), which are those of the WFDB
library (\cref{sec:matlab-library}). The fields are ordinary structure fields, which may
be changed.

\begin{xltabular}{\linewidth}{@{}lL@{}}
\caption{The fields of the result of \code{wfdbx.read\_record}}
\label{tab:matlab-record} \\
\toprule
\tabhead Field & \tabhead Contents \\
\midrule
\endfirsthead
\toprule
\tabhead Field & \tabhead Contents \\
\midrule
\endhead
\bottomrule
\endfoot
\code{name} & the record name, as given to \code{wfdbx.read\_record} \\
\code{fs} & the sampling frequency in Hz, or 0 if the header gives none; with
  \code{'highres', true}, the highest sampling frequency of a multi-frequency record \\
\code{signals} & a matrix of doubles with one row per sample and one column per signal
  read: ADC units, or physical units with \code{'physical', true}; invalid samples are
  \code{wfdbx.WFDB\_INVALID\_SAMPLE()} ($-32768$), or NaN in physical units \\
\code{siginfo} & a 1-by-N \code{wfdbx.siginfo} structure array, one per signal read, in
  the order of the columns; \code{desc} has no white space at either end, and
  \code{adcres} is 0 if the header does not give it \\
\code{names} & a 1-by-N cell array of the signal descriptions, such as \code{'MLII'},
  without the white space that a header may put at either end of them \\
\code{units} & a 1-by-N cell array of the physical units, with \code{[]} where the header
  gives none \\
\code{gains} & a row vector of the gains, in ADC units per physical unit; 0 means that the
  signal is uncalibrated \\
\code{baselines} & a row vector of the ADC values that correspond to 0 physical units \\
\code{base\_time} & the start time and date as text, to the second, such as
  \code{'[14:00:03 21/06/2024]'} (\code{'[14:00:03]'} without a date), or \code{[]} if
  the header gives none \\
\code{base\_time\_ms} & the start time in milliseconds after midnight, such as
  \code{50403500} for 14:00:03.500, or \code{[]} if \code{base\_time} is \code{[]} \\
\code{base\_date} & the start date as \code{[year month day]}, or \code{[]} if the header
  gives none \\
\code{sampfrom} & the number of the first sample returned \\
\code{comments} & a cell column of the info strings of the header (without the
  \code{\#}): its lines that begin with \code{\#} after the record line, found as the
  header parser finds it (an empty or blank first line is not the record line; in legacy
  mode, as the C function \code{getinfo} finds it: the first line that does not begin
  with \code{\#}) \\
\code{checksum\_ok} & \code{false} if the samples do not match the checksums of the
  header; checksums are checked only when a read reaches the end of the record, so a
  shorter read gives \code{true} \\
\end{xltabular}

\begin{xltabular}{\linewidth}{@{}lL@{}}
\caption{The fields of the result of \code{wfdbx.read\_header}}
\label{tab:matlab-header} \\
\toprule
\tabhead Field & \tabhead Contents \\
\midrule
\endfirsthead
\toprule
\tabhead Field & \tabhead Contents \\
\midrule
\endhead
\bottomrule
\endfoot
\code{name} & the record name, as given to \code{wfdbx.read\_header} \\
\code{fs} & the sampling frequency in Hz, or 0 if the header gives none \\
\code{counter\_freq} & the counter frequency, as \code{wfdbx.getcfreq()} returns it (the
  sampling frequency if the header gives none) \\
\code{base\_counter} & the base counter value, as \code{wfdbx.getbasecount()} returns
  it \\
\code{nsamples} & the number of samples per signal, or 0 if the header does not give it \\
\code{nsig} & the number of signals \\
\code{siginfo} & a 1-by-N \code{wfdbx.siginfo} structure array, one per signal;
  \code{desc} has no white space at either end, and \code{adcres} is 0 if the header does
  not give it \\
\code{base\_time}, \code{base\_time\_ms}, \code{base\_date} & the start time and date, as
  for \code{wfdbx.read\_record} \\
\code{comments} & a cell column of the info strings of the header (without the
  \code{\#}), as for \code{wfdbx.read\_record} \\
\code{metadata} & the metadata entries of the header, an N-by-2 cell array
  \code{\{KEY, VALUE; ...\}} in file order, as \code{wfdbx.read\_metadata} returns them:
  from the \file{.hea} file only; \code{cell(0, 2)}, without an error, for an EDF record
  and for a header read from the standard input \\
\end{xltabular}

\begin{xltabular}{\linewidth}{@{}lL@{}}
\caption{The fields of the result of \code{wfdbx.read\_annotations}: one row per
  annotation, in file order}
\label{tab:matlab-annotations} \\
\toprule
\tabhead Field & \tabhead Contents \\
\midrule
\endfirsthead
\toprule
\tabhead Field & \tabhead Contents \\
\midrule
\endhead
\bottomrule
\endfoot
\code{record} & the record name, as given \\
\code{annotator} & the annotator name, such as \code{'atr'} \\
\code{samples} & a column of the annotation times, in samples \\
\code{codes} & a column of the annotation codes, such as \code{wfdbx.NORMAL()} or
  \code{wfdbx.PVC()} \\
\code{symbols} & a cell column of the mnemonics of the codes, such as \code{'N'},
  \code{'V'} or \code{'+'} \\
\code{subtypes} & a column of the subtypes \\
\code{chans} & a column of the channel numbers \\
\code{nums} & a column of the annotator numbers \\
\code{aux} & a cell column of the auxiliary texts (\code{char} vectors), with \code{[]}
  where there is none \\
\code{fs} & the frequency in which \code{samples} are counted, normally the sampling
  frequency of the record; 0 if neither the record, nor the file, nor the option
  \code{'fs'} of \code{wfdbx.read\_annotations} gives one \\
\end{xltabular}

\begin{table}[htbp]
\centering
\caption{The fields of the result of \code{wfdbx.check\_record}}
\label{tab:matlab-checkreport}
\begin{tabularx}{\linewidth}{@{}lL@{}}
\toprule
\tabhead Field & \tabhead Contents \\
\midrule
\code{record} & the record name, as given \\
\code{ok} & \code{true} if no issue has the severity \code{'error'} \\
\code{issues} & an N-by-1 structure array of the issues (\cref{tab:matlab-issue}), in the
  order of the checks \\
\code{lines} & the text form of the report, a cell column: one line per issue,
  \code{'SEVERITY: WHERE: MESSAGE'} (\code{'SEVERITY: MESSAGE'} when \code{where} is
  empty), or the one line \code{'no problems found'} \\
\bottomrule
\end{tabularx}
\end{table}

\begin{table}[htbp]
\centering
\caption{The fields of an issue, a problem found by \code{wfdbx.check\_record}}
\label{tab:matlab-issue}
\begin{tabularx}{\linewidth}{@{}lL@{}}
\toprule
\tabhead Field & \tabhead Contents \\
\midrule
\code{severity} & \code{'error'}, \code{'warning'} or \code{'info'} \\
\code{code} & a stable identifier of the kind of problem, such as \code{'checksum'} or
  \code{'gain\_missing'} (\cref{tab:python-issue-codes}) \\
\code{where} & where the problem is: \code{''} (the record), \code{'signal N'},
  \code{'segment SEG'}, \code{'segment SEG, signal N'} or \code{'annotator ANN'} \\
\code{message} & the description of the problem \\
\bottomrule
\end{tabularx}
\end{table}

\begin{xltabular}{\linewidth}{@{}lL@{}}
\caption{The fields of a \code{wfdbx.siginfo} structure, the specification of a signal}
\label{tab:matlab-siginfo} \\
\toprule
\tabhead Field & \tabhead Contents \\
\midrule
\endfirsthead
\toprule
\tabhead Field & \tabhead Contents \\
\midrule
\endhead
\bottomrule
\endfoot
\code{fname} & the name of the signal file, as the header gives it \\
\code{desc} & the signal description, such as \code{'MLII'}, or \code{[]} \\
\code{units} & the physical units, such as \code{'mV'}, or \code{[]} if not specified \\
\code{gain} & the gain, in ADC units per physical unit (0: uncalibrated) \\
\code{initval} & the value of the first sample \\
\code{group} & the signal group number, from 0 (the signals of one group share one
  file) \\
\code{fmt} & the storage format, such as \code{16}, \code{212} or \code{516} \\
\code{spf} & the number of samples per frame (more than 1 for the faster signals of a
  multi-frequency record) \\
\code{bsize} & the block size, in bytes (0 for ordinary files) \\
\code{adcres} & the ADC resolution, in bits (0: not specified) \\
\code{adczero} & the ADC output for an input of 0\,V \\
\code{baseline} & the ADC value that corresponds to 0 physical units \\
\code{nsamp} & the number of samples (0 if unknown) \\
\code{cksum} & the 16-bit checksum of the samples \\
\end{xltabular}

% ------------------------------------------------------------------------------------------
\subsection{Warnings and errors}
\label{sec:matlab-warnings}

\paragraph{Errors.} The high-level functions report errors with \code{error}, with a
message and an identifier (\cref{tab:matlab-errors}), so that \code{err.identifier} of a
\code{try}/\code{catch} tells which function failed, and, for the metadata functions,
whether an argument was invalid or a file could not be read or written. Each function
checks its arguments before it creates or changes a file, so an invalid argument leaves
the files as they were. After an error, as after a successful call, the settings of the
call (\code{'input\_dir'}, \code{'network'}, \code{'dir'}, \code{'legacy'}), the getvec
mode and the time resolution of output annotators are as they were before it
(\cref{sec:matlab-conventions}). The library may also print a message on the standard error when
it cannot open a file (as in \code{init: can't open header for record REC});
\code{wfdbx.wfdbquiet()} turns these messages off (\cref{sec:matlab-library}).

\begin{table}[htbp]
\centering
\caption{The error and warning identifiers of the high-level API}
\label{tab:matlab-errors}
\begin{tabularx}{\linewidth}{@{}lL@{}}
\toprule
\tabhead Identifier & \tabhead Issued when \\
\midrule
\code{wfdbx:NAME} & an error of the function \code{wfdbx.NAME}: for
  \code{read\_record}, \code{read\_header}, \code{read\_annotations},
  \code{write\_record}, \code{write\_annotations} and \code{check\_record}, any error (a
  file that cannot be found, read, created or written completely; an invalid option or
  argument; a value that the operation needs and the record does not give, such as a
  time in seconds without a sampling frequency); for the five metadata functions, an
  invalid argument: an invalid key or value, a record name that is empty or \code{'-'}, a
  key that the header has more than once (\code{set\_metadata}, \code{update\_metadata}),
  a header without a record line, a header read from a URL to be changed without
  \code{'dir'}, or an invalid option \\
\code{wfdbx:NAME:file} & a metadata function cannot read or write a file: a header that
  cannot be found or read, a changed header that cannot be written, or a \code{'dir'}
  that does not exist \\
\code{wfdbx:checksum} & (a warning) the samples read by \code{wfdbx.read\_record} do not
  match the checksums of the header \\
\code{wfdbx:duplicate\_key} & (a warning) a header has more than one metadata entry for a
  key \\
\bottomrule
\end{tabularx}
\end{table}

\paragraph{The checksum warning.} When \code{wfdbx.read\_record} reads to the end of a
record whose samples do not match the checksums of its header, it returns all the samples
and issues the warning \code{wfdbx:checksum}: ``checksum error in record REC: the signal
data do not match the checksums in its header''. \code{'checksum\_warning', false} turns
the warning off; \code{rec.checksum\_ok} is \code{false} either way. This applies only
when every sample could be read: a signal file that ends before the length that its
header gives is an error (``can't read signal data of record REC (getvec status -3)''),
even though the checksum, verified at the frame in which the file ended, does not match
either; \code{wfdbx.check\_record} reports it as \code{length\_short}, not as a
checksum error.

\paragraph{The duplicate-key warning.} When a header holds two metadata entries whose keys
differ only in letter case (written by other software), \code{wfdbx.read\_metadata},
\code{wfdbx.get\_metadata} and \code{wfdbx.read\_header} use the first and issue the
warning \code{wfdbx:duplicate\_key}: ``header of record REC has more than one entry for
key KEY (lines L1, L2); the first is used''. \code{wfdbx.set\_metadata} and
\code{wfdbx.update\_metadata} then refuse to change that key (an argument error), and
\code{wfdbx.remove\_metadata} removes all of its entries.

Both are issued with \code{warning}, so \code{warning('off', 'wfdbx:checksum')} turns one
off, \code{lastwarn} returns the last one, and \code{warning('error', ...)} (MATLAB) makes
it an error, as for any other warning:

\begin{lstlisting}[language=Matlab,style=example]
% A record whose last sample is changed after it was written: the checksum fails.
wfdbx.write_record('ck', (0:99)', 100, 'gains', 100, 'adcres', 12, 'units', 'mV');
fid = fopen('ck.dat', 'r+b');
fseek(fid, -2, 'eof');
fwrite(fid, [0 0], 'uint8');
fclose(fid);
lastwarn('');
rec = wfdbx.read_record('ck');                  % prints the checksum warning
fprintf('%d %d %d\n', size(rec.signals), rec.checksum_ok);   % 100 1 0: all the samples
[msg, id] = lastwarn();
disp(id);                                       % wfdbx:checksum
rec = wfdbx.read_record('ck', 'checksum_warning', false);    % no warning

% Two entries for one key: the first is used, with a warning.
wfdbx.write_record('dup', zeros(10, 1), 100, 'gains', 0, 'adcres', 12, ...
    'comments', {' Age: 63', ' AGE: 64'});
disp(wfdbx.get_metadata('dup', 'age'));         % 63, with the warning
[msg, id] = lastwarn();
disp(id);                                       % wfdbx:duplicate_key
try
    wfdbx.set_metadata('dup', 'Age', '65');
catch err
    fprintf('%s: %s\n', err.identifier, err.message);
    % wfdbx:set_metadata: header of record dup has more than one entry for key Age
    % (lines 3, 4); remove them first
end
disp(wfdbx.remove_metadata('dup', 'age'));      % 2
\end{lstlisting}

\begin{lstlisting}[language=Matlab,style=example]
try
    wfdbx.read_record('nosuchrecord');
catch err
    fprintf('%s: %s\n', err.identifier, err.message);
    % wfdbx:read_record: can't read record nosuchrecord
end
try
    wfdbx.read_metadata('nosuchrecord');
catch err
    fprintf('%s: %s\n', err.identifier, err.message);
    % wfdbx:read_metadata:file: can't read header for record nosuchrecord
end
wfdbx.write_record('raw', (0:9)', 100, 'gains', 0, 'adcres', 12);   % uncalibrated
try
    wfdbx.read_record('raw', 'physical', true);
catch err
    fprintf('%s: %s\n', err.identifier, err.message);
    % wfdbx:read_record: signal 0 of record raw is uncalibrated
end
\end{lstlisting}

% ==========================================================================================
\section{The WFDB library functions in MATLAB and Octave}
\label{sec:matlab-library}

Every function of the WFDB library is a function of the package \code{wfdbx} under its C
name, with the same arguments, in the same order, and the same return values, so that
programs and documentation written for the C library carry over directly. The functions,
their arguments and their results are documented once, for all the languages, in
\cref{chap:library}; this section gives the rules by which the C declarations become
MATLAB functions (see also \cref{sec:library-matlab}). \code{help wfdbx.NAME} gives the
documentation of each function, including what WFDB-x changes in it.

\subsection{Conventions}

\paragraph{Names.} The functions keep their C names, with the prefix of the package:
\code{wfdbx.isigopen}, \code{wfdbx.getvec}, \ldots{} The constants (\code{WFDB\_READ},
\code{WFDB\_INVALID\_SAMPLE}, \code{WFDB\_HIGHRES}, \ldots) and the annotation codes
(\code{NORMAL}, \code{PVC}, \ldots, \cref{chap:formats}) are functions without arguments:
\code{wfdbx.WFDB\_READ()}, \code{wfdbx.NORMAL()}; \code{C = wfdbx.ecgcodes()} returns all
the annotation codes as the fields of one structure (\code{C.NORMAL}, \code{C.PVC},
\ldots), which is faster in loops. The calibration constants (\code{WFDB\_AC\_COUPLED},
\ldots) and \code{WFDB\_FMT\_LIST} are not defined; their values are given in
\cref{tab:library-constants}. The macros of \file{<wfdb/ecgmap.h>} are functions, under
both their macro names and the names of the functions behind them:
\code{wfdbx.isqrs} and \code{wfdbx.wfdb\_isqrs}, \code{wfdbx.map1} and
\code{wfdbx.wfdb\_map1}, and so on.

\paragraph{NULL is \code{[]}.} A C \code{NULL} string, as argument or result, is
\code{[]} (an empty double), and an empty C string is \code{''}; \code{ischar} tells them
apart. \code{wfdbx.sampfreq([])} returns the frequency of the open record,
\code{wfdbx.getinfo([])} continues with the info strings of the last record, and
\code{wfdbx.getinfo} returns \code{[]} after the last one.

\paragraph{Arrays and output arguments.} Where a C function fills an array or a structure
that the caller passes, the MATLAB function returns it as an additional output value, and
omits it from the arguments:
\begin{lstlisting}[language=Matlab]
[nsig, si] = wfdbx.isigopen(record, nsig)     % si: a struct array of siginfo
[nsig, si] = wfdbx.osigopen(record, nsig)
[stat, si] = wfdbx.wfdbinit(record, aiarray, nann, nsig)
[stat, v] = wfdbx.getvec(v0)                  % v0 (optional): the vector C leaves unchanged
[stat, v] = wfdbx.getframe(v0)
[stat, annot] = wfdbx.getann(an, annot0)      % annot0 (optional), as v0
[stat, cal] = wfdbx.getcal(desc, units)
[stat, si] = wfdbx.setheader(record, si, nsig)
segs = wfdbx.getseginfo()                     % a struct array of seginfo
[r, argv] = wfdbx.wfdbx_legacy_args(argv)     % argv: a cell array
\end{lstlisting}
With \code{nsig} $\le 0$, \code{wfdbx.isigopen} reads the header only, and \code{si}
describes \code{-nsig} signals. Arrays of structures passed in are struct arrays:
\code{wfdbx.annopen(record, wfdbx.anninfo('atr', wfdbx.WFDB\_READ()))},
\code{wfdbx.osigfopen(si, nsig)}. \code{wfdbx.setmsheader(record, segnames)} takes a cell
array of segment names.

\paragraph{Numbers and strings.} Samples, times, counts and frequencies are doubles,
exact for every 32-bit sample and for times up to $2^{53}$. Strings are \code{char} row
vectors whose characters are the bytes of the C strings (codes 0 to 255); record names
may also be MATLAB strings. The \code{aux} field of an annotation is \code{[]}, or a
\code{uint8} row vector whose first element is the length of the text that follows, as in
C: \code{char(a.aux(2:1 + a.aux(1)))} is the text.

\paragraph{Return values and errors.} The functions return what the C functions return: a
status (negative on errors), a number of signals, a frequency, a string or \code{[]}.
Errors are reported by these values and by a message on the standard error, as in C,
never by MATLAB errors. \code{wfdbx.wfdberror()} returns the last message;
\code{wfdbx.wfdbquiet()} turns the messages off and \code{wfdbx.wfdbverbose()} on again;
\code{wfdbx.wfdb\_error(fmt, ...)} reports a message as the library does, formatting as
C's \code{printf}.

\paragraph{State.} As in C, the library keeps its state in the session (in
\code{persistent} variables): the open record, its input signals, the output signals and
the annotators remain open between calls, until \code{wfdbx.wfdbquit()} closes them.
\code{wfdbx.wfdbquit()} leaves the input and output directories, the network setting,
legacy mode, the getvec mode and the time resolution of output annotators (which opening
an EDF record sets, as in C) as they are; \code{clear functions} (or \code{clear all}) resets the whole
state. The high-level functions of \cref{sec:matlab-api} (except the metadata functions)
call \code{wfdbx.wfdbquit()}, which closes anything opened with the library functions.

\paragraph{Additions.} \code{[v, stat, last] = wfdbx.getframes(n, minstat)} reads up to
\code{n} sample vectors at once (\code{n} may be \code{Inf}), as \code{n} calls of
\code{wfdbx.getvec} would, but decoding whole blocks of the signal files: \code{v} is a
matrix with one row per vector and one column per signal, \code{stat} the status of the
first vector that could not be read, or the number of signals if all \code{n} were read,
and \code{last} the vector filled by the call that stopped the reading
(\cref{sec:library-additions}). \code{stat = wfdbx.putframes(M)} writes the rows of
\code{M}, as \code{wfdbx.putvec} for each row would. Calling \code{wfdbx.getvec} sample by
sample is much slower in MATLAB and Octave than these functions, and than the high-level
functions.

\subsection{The structures}
\label{sec:matlab-structures}

The structures of the C library are MATLAB structures with the names of the fields of
the C structures (\cref{tab:matlab-structures}), made with C's zero values by these
functions:
\begin{lstlisting}[language=Matlab]
si = wfdbx.siginfo(n)            % 1-by-n struct array (n = 1 by default); spf is 1
s = wfdbx.seginfo(n)             % 1-by-n struct array (n = 1 by default)
ai = wfdbx.anninfo(name, stat)   % one structure
a = wfdbx.annotation(n)          % 1-by-n struct array (n = 1 by default)
cal = wfdbx.calinfo()            % one structure
\end{lstlisting}

\begin{xltabular}{\linewidth}{@{}lL@{}}
\caption{The C structures in MATLAB and Octave}
\label{tab:matlab-structures} \\
\toprule
\tabhead Constructor (C structure) & \tabhead Fields \\
\midrule
\endfirsthead
\toprule
\tabhead Constructor (C structure) & \tabhead Fields \\
\midrule
\endhead
\bottomrule
\endfoot
\code{wfdbx.siginfo} (\code{WFDB\_Siginfo}) & \code{fname}, \code{desc}, \code{units}
  (\code{[]} or \code{char}), \code{gain}, \code{initval}, \code{group}, \code{fmt},
  \code{spf}, \code{bsize}, \code{adcres}, \code{adczero}, \code{baseline}, \code{nsamp},
  \code{cksum} (\cref{tab:matlab-siginfo}) \\
\code{wfdbx.seginfo} (\code{WFDB\_Seginfo}) & \code{recname} (the name of the segment),
  \code{nsamp} (its number of samples), \code{samp0} (the number of its first sample in
  the record) \\
\code{wfdbx.anninfo} (\code{WFDB\_Anninfo}) & \code{name} (the annotator name, the suffix
  of the file), \code{stat} (\code{wfdbx.WFDB\_READ()} or \code{wfdbx.WFDB\_WRITE()} for
  the MIT format, \code{wfdbx.WFDB\_AHA\_READ()} or \code{wfdbx.WFDB\_AHA\_WRITE()} for the
  AHA format) \\
\code{wfdbx.annotation} (\code{WFDB\_Annotation}) & \code{time} (in sample intervals),
  \code{anntyp} (the annotation code), \code{subtyp}, \code{chan}, \code{num}, \code{aux}
  (\code{[]}, or a \code{uint8} row vector whose first element is the length of the text
  that follows) \\
\code{wfdbx.calinfo} (\code{WFDB\_Calinfo}) & \code{low}, \code{high} (the levels of the
  calibration pulse, in physical units), \code{scale} (physical units per cm),
  \code{sigtype} (matched against the beginning of a signal description), \code{units},
  \code{caltype} \\
\end{xltabular}

\subsection{The settings functions}
\label{sec:matlab-settings}

The functions of \cref{tab:matlab-settings} are additions of WFDB-x to the library: they
set where files are read and written, network access and legacy mode, for the high-level
functions, the library functions and the tools alike (\cref{chap:files,chap:behavior}),
and give the path under which a program opens a file that it reads or creates itself, as
the tools do (\code{wfdbx.wfdbx\_input\_path}, \code{wfdbx.wfdbx\_output\_path};
\cref{sec:tool-files}). Each \code{set} function returns the previous setting, so that it
can be restored. The C function \code{wfdbx.setwfdb(p)} sets the input directory too (\code{p} must name at most
one directory), and \code{wfdbx.getwfdb()} returns it; in legacy mode they set and return
the WFDB path, as in C.

\begin{xltabular}{\linewidth}{@{}lL@{}}
\caption{The settings functions}
\label{tab:matlab-settings} \\
\toprule
\tabhead Function & \tabhead Effect and result \\
\midrule
\endfirsthead
\toprule
\tabhead Function & \tabhead Effect and result \\
\midrule
\endhead
\bottomrule
\endfoot
\code{previous = wfdbx.wfdbx\_set\_input\_dir(d)} & sets the input directory: a local
  folder, a \code{\textasciitilde} path or a URL; \code{[]} or \code{''} restores the
  default (\code{WFDBX\_INPUT\_DIR}, or else the current folder). Returns the previous
  input directory, or \code{[]} if none was set \\
\code{d = wfdbx.wfdbx\_get\_input\_dir()} & returns the input directory, or \code{[]} if
  none is set \\
\code{previous = wfdbx.wfdbx\_set\_output\_dir(d)} & sets the output directory: a local
  folder or a \code{\textasciitilde} path (a URL is refused, with a message, and the
  setting is unchanged); \code{[]} or \code{''} restores the default
  (\code{WFDBX\_OUTPUT\_DIR}, or else the current folder). Returns the previous output
  directory, or \code{[]} \\
\code{d = wfdbx.wfdbx\_get\_output\_dir()} & returns the output directory, or \code{[]}
  if none is set \\
\code{previous = wfdbx.wfdbx\_set\_network(on)} & turns network access on (\code{true} or
  1) or off (\code{false} or 0); \code{[]} or a negative number restores the default (on
  if \code{WFDBX\_NETWORK} is a nonzero number). Returns the previous setting, 1 or 0 \\
\code{on = wfdbx.wfdbx\_get\_network()} & returns 1 if network access is on, 0
  otherwise \\
\code{previous = wfdbx.wfdbx\_set\_legacy(on)} & turns legacy mode on (\code{true} or 1)
  or off, for the whole library. Returns the previous setting, 1 or 0 \\
\code{on = wfdbx.wfdbx\_get\_legacy()} & returns 1 if legacy mode is on, 0 otherwise \\
\code{[r, argv] = wfdbx.wfdbx\_legacy\_args(argv)} & if \code{argv\{2\}} (\code{argv}: a
  cell row of the arguments, program name first) is \code{'-{}-legacy'}, removes it,
  turns legacy mode on and returns \code{r} = 1; otherwise returns \code{argv} unchanged
  and \code{r} = 0 \\
\code{r = wfdbx.wfdbx\_time\_needs\_freq(string)} & returns 1 if
  \code{wfdbx.strtim(string)} needs a sampling frequency, 0 otherwise \\
\code{n = wfdbx.wfdbx\_write\_errors()} & returns the number of write errors since the
  program started (output files that could not be written, flushed or closed) \\
\code{wfdbx.wfdbx\_note\_write\_error()} & counts a write error, for a program that
  writes a file itself \\
\code{p = wfdbx.wfdbx\_input\_path(name)} & returns the path under which a program opens
  a file that it reads itself: in normal mode, a relative name under the input directory,
  a \code{\textasciitilde} name expanded, an absolute name as given; in legacy mode, the
  name as given. \code{[]}, \code{''} and \code{'-'} are returned as given; \code{[]} for
  a URL, or a name that becomes one under the input directory \\
\code{p = wfdbx.wfdbx\_output\_path(name)} & returns the path under which a program
  creates a file that it writes itself: in normal mode, a relative name under the output
  directory (whose missing folders are then created), a \code{\textasciitilde} name
  expanded, an absolute name as given; in legacy mode, any name under the output
  directory, as a plain prefix. \code{[]}, \code{''} and \code{'-'} are returned as given;
  \code{[]} for a URL \\
\end{xltabular}

\subsection{Examples}

Reading the samples and the annotations of a record, as a C program does (the record
\code{100s} is read from the current folder):

\begin{lstlisting}[language=Matlab,style=example]
[nsig, si] = wfdbx.isigopen('100s', 2);         % open 2 signals
fprintf('%d %s %d %g\n', nsig, si(1).desc, si(1).fmt, si(1).gain);   % 2 MLII 212 200
fprintf('%g\n', wfdbx.sampfreq([]));            % 360: the frequency of the open record

wfdbx.isigsettime(wfdbx.strtim('0:30'));        % skip to 30 seconds
[stat, v] = wfdbx.getvec();                     % one sample per signal
fprintf('%d %d %g\n', v(1), v(2), wfdbx.aduphys(0, v(1)));      % 947 968 -0.385

[frames, stat] = wfdbx.getframes(360);          % the next second, as a matrix
fprintf('%d %d %d\n', size(frames), stat);      % 360 2 2

wfdbx.annopen('100s', wfdbx.anninfo('atr', wfdbx.WFDB_READ()));
n = 0;
[stat, a] = wfdbx.getann(0);
while stat == 0                                 % 0: an annotation was read; -1: the end
    if n < 2
        fprintf('%s %s\n', wfdbx.mstimstr(a.time), wfdbx.annstr(a.anntyp));
    end
    n = n + 1;
    [stat, a] = wfdbx.getann(0);
end
fprintf('%d annotations\n', n);                 % 75 annotations

wfdbx.wfdbquiet();                              % no messages on the standard error
stat = wfdbx.isigopen('nosuchrecord', 0);
fprintf('%d %s', stat, wfdbx.wfdberror());
% -1 init: can't open header for record nosuchrecord
wfdbx.wfdbverbose();
wfdbx.wfdbquit();                               % close everything
\end{lstlisting}

Writing a record and an annotation file, sample by sample and annotation by annotation:

\begin{lstlisting}[language=Matlab,style=example]
out = wfdbx.siginfo(1);
out.fname = 'ramp.dat';
out.desc = 'ramp';
out.units = 'mV';
out.gain = 200;
out.fmt = 16;
out.adcres = 12;
disp(wfdbx.osigfopen(out, 1));      % 1: one output signal
wfdbx.setsampfreq(250);
for k = 0:499
    wfdbx.putvec(k);
end
disp(wfdbx.newheader('ramp'));      % 0: ramp.hea is written
wfdbx.annopen('ramp', wfdbx.anninfo('qrs', wfdbx.WFDB_WRITE()));
a = wfdbx.annotation();
a.anntyp = wfdbx.NORMAL();
for t = [100 350]
    a.time = t;
    wfdbx.putann(0, a);
end
wfdbx.wfdbquit();                   % completes ramp.dat and ramp.qrs

ann = wfdbx.read_annotations('ramp', 'qrs');
disp(ann.samples');                 % 100 350
\end{lstlisting}

% ==========================================================================================
\section{The command-line tools in MATLAB and Octave}
\label{sec:matlab-tools}

The 33 command-line tools (\code{rdsamp}, \code{rdann}, \code{wrsamp}, \code{xform},
\ldots) are functions of the package \code{wfdbx.apps}, one file per tool in
\file{+wfdbx/+apps}. The tools, their options, their output and their exit status are
those of the C tools, and are documented in \cref{chap:tools}; \code{help
wfdbx.apps.TOOL} shows the usage of a tool.

\paragraph{From the shell, with Octave.} The Octave script \file{matlab/bin/wfdbx\_cli.m}
runs a tool: its first argument names the tool, and the others are the tool's arguments.
It adds \file{matlab} to the path, and exits with the tool's status. An alias makes the
command shorter:

\begin{lstlisting}[style=shell]
alias wfdbx='octave-cli --no-gui --norc -q /path/to/wfdb-x/matlab/bin/wfdbx_cli.m'
export WFDBX_INPUT_DIR=/path/to/wfdb-x/tests/data
wfdbx rdsamp -r 100s -t 10 -p     # the first 10 seconds of 100s, in physical units
wfdbx rdann -r 100s -a atr        # its annotations
wfdbx rdsamp --legacy -r 100s     # legacy mode: --legacy as the first argument
\end{lstlisting}

The script also takes these arguments, instead of a tool:

\begin{table}[htbp]
\centering
\caption{The arguments of the script \code{wfdbx\_cli.m} (as \cmd{wfdbx})}
\label{tab:matlab-wfdbx-options}
\begin{tabularx}{\linewidth}{@{}lL@{}}
\toprule
\tabhead Command & \tabhead Effect \\
\midrule
\cmd{wfdbx TOOL [ARGUMENTS ...]} & runs the tool \code{TOOL}, and exits with its
  status \\
\cmd{wfdbx -{}-list} & lists the names of the tools, one per line \\
\cmd{wfdbx -{}-version} & prints the version, \code{WFDB-x 1.2.2} \\
\cmd{wfdbx -{}-help}, \cmd{wfdbx -h} & prints the usage \\
\bottomrule
\end{tabularx}
\end{table}

Without arguments, the script prints the usage and exits with status 1; an unknown tool
is an error (``wfdbx: unknown tool 'TOOL' (try 'wfdbx -{}-list')'', status 1), and a
MATLAB error in a tool is reported as \code{wfdbx: error: MESSAGE} (status 1). When an
output file cannot be written completely (a full disk, for example), or the standard
output cannot be written (``TOOL: write error on standard output''), a tool exits with
status 1; in legacy mode, it keeps the original status, as in C. Octave starts in a fraction of a second, so
the script suits everyday use.

\paragraph{From the shell, with MATLAB.} The shell script
\file{matlab/bin/wfdbx\_matlab.sh} runs a tool with \cmd{matlab -batch}, or the program
named by the environment variable \code{WFDBX\_MATLAB}. It saves its standard input to a
temporary file (named by the environment variable \code{WFDBX\_STDIN}) before it starts
MATLAB, so a script whose standard input stays open runs it with \cmd{< /dev/null}, and it
passes the arguments in hexadecimal, so that any bytes arrive unchanged. MATLAB starts
slowly, so the script is mainly useful for testing.

\begin{lstlisting}[style=shell]
export WFDBX_INPUT_DIR=/path/to/wfdb-x/tests/data
/path/to/wfdb-x/matlab/bin/wfdbx_matlab.sh rdsamp -r 100s -t 10 -p
\end{lstlisting}

\paragraph{In a session.} Each tool is a function that takes the arguments as a cell
array (without the program name), writes to the standard output and error, and returns
the exit status:
\begin{lstlisting}[language=Matlab]
status = wfdbx.apps.rdsamp({'-r', '100s', '-t', 's2'})    % the first 2 samples
\end{lstlisting}
A call of \code{exit} in a C tool ends the function, not the session.

\paragraph{Environment variables.} The tools, like the rest of the package, use the
environment variables of \cref{tab:matlab-environment}, also for the files that they read
or create themselves (\cref{sec:tool-files}): the \opt{-i} file of \cmd{wrsamp},
\cmd{edf2mit}, \cmd{wav2mit}, \cmd{a2m} and \cmd{ad2m}, the input files of
\cmd{ahaecg2mit} and the \opt{-S} script of \cmd{xform} are read from the input directory,
and the WAV file of \cmd{mit2wav}, the EDF file of \cmd{mit2edf}, the master header of
\cmd{snip -m} and the \file{.des}, \file{.ecg}, \file{.key} and \file{.txt} files of
\cmd{parsescp} are created in the output directory. Such a file cannot be read from a URL.
In legacy mode, the files that a tool reads itself are opened as named, and the output
directory, if set, is a plain prefix of every output name. In Octave and in MATLAB,
\code{getenv} returns \code{''} for an unset and an empty variable alike, so an empty
variable counts as unset.

\begin{table}[htbp]
\centering
\caption{The environment variables used by the MATLAB/Octave package}
\label{tab:matlab-environment}
\begin{tabularx}{\linewidth}{@{}lL@{}}
\toprule
\tabhead Variable & \tabhead Meaning \\
\midrule
\code{WFDBX\_INPUT\_DIR} & the input directory (a folder, a \code{\textasciitilde} path or
  a URL), when \code{wfdbx.wfdbx\_set\_input\_dir} does not set one \\
\code{WFDBX\_OUTPUT\_DIR} & the output directory, when \code{wfdbx.wfdbx\_set\_output\_dir}
  does not set one \\
\code{WFDBX\_NETWORK} & network access is on if it is a nonzero number, when
  \code{wfdbx.wfdbx\_set\_network} does not set it \\
\code{WFDBANNSORT} & \code{0} to keep annotation files in the order written (no sorting
  when they are closed) \\
\code{WFDBGVMODE} & the default getvec mode (\code{setgvmode} with a negative mode) \\
\code{WFDBCAL} & the calibration file read by \code{wfdbx.calopen([])} (default
  \file{wfdbcal}) \\
\code{WFDB} & the WFDB path, used in legacy mode only (outside legacy mode, it is
  ignored, with a message) \\
\code{WFDB\_FLAC\_CHECK\_MD5} & if set, the MD5 signature of FLAC signal files is
  verified when they are read \\
\code{WFDBX\_MATLAB} & the MATLAB command that \file{wfdbx\_matlab.sh} runs (default
  \cmd{matlab}) \\
\bottomrule
\end{tabularx}
\end{table}

Tools that write signal parameters require them as options: \code{wrsamp} (with
\opt{-o}) requires \opt{-F} (frequency), \opt{-G} (gain), \opt{-R} (ADC resolution) and,
for calibrated signals, \opt{-U} (units); \code{wav2mit} requires \opt{-G}, and \opt{-U}
unless the gain is 0; \code{ad2m} and \code{ahaecg2mit} require \opt{-F}, \opt{-G},
\opt{-R}, and \opt{-U} unless the gain is 0; \code{a2m}, \code{rdann}, \code{ann2rr},
\code{sortann} and \code{m2a} require \opt{-F} when the record's header gives no
frequency (\cref{chap:tools}).

% ==========================================================================================
\section{Testing}
\label{sec:matlab-testing}

The unit tests of the MATLAB/Octave implementation are in \file{matlab/tests}: one
function \file{test\_*.m} per group, run by \code{run\_tests}. From the root of the
repository:

\begin{lstlisting}[style=shell]
octave-cli --no-gui --norc -q --eval "cd matlab/tests; run_tests"           # Octave
matlab -batch "cd matlab/tests; run_tests"                                  # MATLAB
octave-cli --no-gui --norc -q --eval "cd matlab/tests; run_tests('metadata')"  # by name
\end{lstlisting}

\code{run\_tests} prints \code{ok NAME} or \code{FAIL NAME: MESSAGE} for each test, then
a summary, and ends with an error (so that the command exits with status 1) if a test
failed; \code{run\_tests(PATTERN)} runs the tests whose names match a regular expression,
and \code{n = run\_tests(...)} returns the number of failures. The tests cover the C
standard library emulation, the WFDB library functions (with expected values printed by
the C library), the high-level functions, the bulk readers and writers against the
per-sample functions, the FLAC codec, finding files, \code{wfdbx.check\_record}, and the
metadata functions, with the cases shared by the seven languages (in
\file{tests/metadata}: results, warnings, and the headers, byte for byte). They run in
temporary folders, with copies of the test records of \file{tests/data}. The tests of
network access use a local HTTP server on 127.0.0.1 only (a Python script, started with
\cmd{python3}; without it, they are skipped).

The cross-language conformance suite runs every tool through \file{matlab/bin/wfdbx\_cli.m}
and compares its results, byte for byte, with those of the original C programs (1155
cases), and, with \opt{-{}-legacy}, in legacy mode (886 cases). From the root of the
repository:

\begin{lstlisting}[style=shell]
PYTHONPATH=$PWD/python/src python3 tests/conformance/harness.py check --impl octave -j8
PYTHONPATH=$PWD/python/src python3 tests/conformance/harness.py check --impl octave -j8 --legacy
\end{lstlisting}

The implementation \code{matlab} of the harness runs the tools through
\file{matlab/bin/wfdbx\_matlab.sh} instead.
